\appendix

\chapter{The Interdisciplinary Translation Matrix}
\label{app:translation_matrix}

Because the AVE framework roots physical reality in the deterministic continuum mechanics of a discrete substrate graph, its foundational equations project symmetrically outward into multiple established disciplines of applied engineering and mathematics. The framework serves as a universal translation matrix between abstract Quantum Field Theory (QFT) and classical macroscopic disciplines.

\section{The Rosetta Stone of Physics}

The following domain-specific translation tables are each defined once in a canonical location (\texttt{manuscript/common/}) and \texttt{\textbackslash input}'d into the relevant chapter.  The complete set is collected here as a single reference.

\subsection{Circuit Analysis (Topo-Kinematic Identity)}
% ═══════════════════════════════════════════════════════════════
% DOMAIN TRANSLATION TABLE: Topo-Kinematic Circuit Identity
% Source: Extracted from Book 5 §12
% Canonical location: manuscript/common/translation_circuit.tex
% ═══════════════════════════════════════════════════════════════
%
% Usage: % ═══════════════════════════════════════════════════════════════
% DOMAIN TRANSLATION TABLE: Topo-Kinematic Circuit Identity
% Source: Extracted from Book 5 §12
% Canonical location: manuscript/common/translation_circuit.tex
% ═══════════════════════════════════════════════════════════════
%
% Usage: % ═══════════════════════════════════════════════════════════════
% DOMAIN TRANSLATION TABLE: Topo-Kinematic Circuit Identity
% Source: Extracted from Book 5 §12
% Canonical location: manuscript/common/translation_circuit.tex
% ═══════════════════════════════════════════════════════════════
%
% Usage: \input{../../common/translation_circuit.tex}
%        (from any chapter two levels deep in manuscript/)

\begin{table}[H]
\centering
\caption{Topo-Kinematic Circuit Identity.  Every row is derived from $\xi_{topo} = e/\ell_{node} \approx 4.149 \times 10^{-7}$ C/m (Axioms~1--2).}
\label{tab:trans_circuit}
\small
\renewcommand{\arraystretch}{1.3}
\begin{tabular}{@{}l l c l@{}}
\toprule
\textbf{Electrical} & \textbf{Mechanical} & \textbf{Mapping} & \textbf{Units Check} \\
\midrule
Charge $Q$        & Displacement $x$ & $Q = \xi\, x$             & $[\text{C/m}][\text{m}] = [\text{C}]$ \\
Current $I$       & Velocity $v$     & $I = \xi\, v$             & $[\text{C/m}][\text{m/s}] = [\text{A}]$ \\
Voltage $V$       & Force $F$        & $V = \xi^{-1} F$          & $[\text{m/C}][\text{N}] = [\text{V}]$ \\
Inductance $L$    & Mass $m$         & $L = \xi^{-2} m$          & $[\text{m}^2/\text{C}^2][\text{kg}] = [\text{H}]$ \\
Capacitance $C$   & Compliance $\kappa$ & $C = \xi^{2} \kappa$   & $[\text{C}^2/\text{m}^2][\text{m/N}] = [\text{F}]$ \\
Resistance $R$    & Viscosity $\eta$ & $R = \xi^{-2} \eta$       & $[\text{m}^2/\text{C}^2][\text{kg/s}] = [\Omega]$ \\
\bottomrule
\end{tabular}
\end{table}

%        (from any chapter two levels deep in manuscript/)

\begin{table}[H]
\centering
\caption{Topo-Kinematic Circuit Identity.  Every row is derived from $\xi_{topo} = e/\ell_{node} \approx 4.149 \times 10^{-7}$ C/m (Axioms~1--2).}
\label{tab:trans_circuit}
\small
\renewcommand{\arraystretch}{1.3}
\begin{tabular}{@{}l l c l@{}}
\toprule
\textbf{Electrical} & \textbf{Mechanical} & \textbf{Mapping} & \textbf{Units Check} \\
\midrule
Charge $Q$        & Displacement $x$ & $Q = \xi\, x$             & $[\text{C/m}][\text{m}] = [\text{C}]$ \\
Current $I$       & Velocity $v$     & $I = \xi\, v$             & $[\text{C/m}][\text{m/s}] = [\text{A}]$ \\
Voltage $V$       & Force $F$        & $V = \xi^{-1} F$          & $[\text{m/C}][\text{N}] = [\text{V}]$ \\
Inductance $L$    & Mass $m$         & $L = \xi^{-2} m$          & $[\text{m}^2/\text{C}^2][\text{kg}] = [\text{H}]$ \\
Capacitance $C$   & Compliance $\kappa$ & $C = \xi^{2} \kappa$   & $[\text{C}^2/\text{m}^2][\text{m/N}] = [\text{F}]$ \\
Resistance $R$    & Viscosity $\eta$ & $R = \xi^{-2} \eta$       & $[\text{m}^2/\text{C}^2][\text{kg/s}] = [\Omega]$ \\
\bottomrule
\end{tabular}
\end{table}

%        (from any chapter two levels deep in manuscript/)

\begin{table}[H]
\centering
\caption{Topo-Kinematic Circuit Identity.  Every row is derived from $\xi_{topo} = e/\ell_{node} \approx 4.149 \times 10^{-7}$ C/m (Axioms~1--2).}
\label{tab:trans_circuit}
\small
\renewcommand{\arraystretch}{1.3}
\begin{tabular}{@{}l l c l@{}}
\toprule
\textbf{Electrical} & \textbf{Mechanical} & \textbf{Mapping} & \textbf{Units Check} \\
\midrule
Charge $Q$        & Displacement $x$ & $Q = \xi\, x$             & $[\text{C/m}][\text{m}] = [\text{C}]$ \\
Current $I$       & Velocity $v$     & $I = \xi\, v$             & $[\text{C/m}][\text{m/s}] = [\text{A}]$ \\
Voltage $V$       & Force $F$        & $V = \xi^{-1} F$          & $[\text{m/C}][\text{N}] = [\text{V}]$ \\
Inductance $L$    & Mass $m$         & $L = \xi^{-2} m$          & $[\text{m}^2/\text{C}^2][\text{kg}] = [\text{H}]$ \\
Capacitance $C$   & Compliance $\kappa$ & $C = \xi^{2} \kappa$   & $[\text{C}^2/\text{m}^2][\text{m/N}] = [\text{F}]$ \\
Resistance $R$    & Viscosity $\eta$ & $R = \xi^{-2} \eta$       & $[\text{m}^2/\text{C}^2][\text{kg/s}] = [\Omega]$ \\
\bottomrule
\end{tabular}
\end{table}


\subsection{Quantum Mechanics}
% ═══════════════════════════════════════════════════════════════
% DOMAIN TRANSLATION TABLE: Quantum Mechanics ↔ AVE
% Source: Extracted from Book 3 Ch.16
% Canonical location: manuscript/common/translation_qm.tex
% ═══════════════════════════════════════════════════════════════
%
% Usage: % ═══════════════════════════════════════════════════════════════
% DOMAIN TRANSLATION TABLE: Quantum Mechanics ↔ AVE
% Source: Extracted from Book 3 Ch.16
% Canonical location: manuscript/common/translation_qm.tex
% ═══════════════════════════════════════════════════════════════
%
% Usage: % ═══════════════════════════════════════════════════════════════
% DOMAIN TRANSLATION TABLE: Quantum Mechanics ↔ AVE
% Source: Extracted from Book 3 Ch.16
% Canonical location: manuscript/common/translation_qm.tex
% ═══════════════════════════════════════════════════════════════
%
% Usage: \input{../../common/translation_qm.tex}

\begin{table}[H]
\centering
\caption{Quantum Mechanics $\leftrightarrow$ AVE Translation.  Rows 1--8: same mathematics, different ontology.  Row~9: genuinely new AVE prediction.}
\label{tab:trans_qm}
\small
\renewcommand{\arraystretch}{1.3}
\begin{tabular}{@{}p{3.2cm} p{3.5cm} p{7cm}@{}}
\toprule
\textbf{Standard QM} & \textbf{AVE Equivalent} & \textbf{Relationship} \\
\midrule
Wavefunction $\psi(r)$ & Acoustic cavity mode & Same equation (Helmholtz). AVE: mode of vacuum LC mesh. \\
$|\psi|^2$ probability & Time-averaged trajectory & Density of point-defect sweeping its standing-wave mode. \\
Coulomb potential $e^2/(4\pi\varepsilon_0 r)$ & $\alpha\hbar c / r$ (Axiom 2) & Algebraically identical. $K = \alpha\hbar c$ is the impedance coupling. \\
Bohr radius $a_0$ & $\ell_{node}/\alpha$ & $a_0 = \hbar/(m_e c \alpha)$. The atom is $1/\alpha \approx 137$ lattice pitches. \\
Hartree energy $E_H$ & $m_e(\alpha c)^2$ & Same formula. $E_H = K/a_0$ (coupling per cavity radius). \\
Electron & $0_1$ topological unknot & Irreducible closed flux tube. Internal radius $\sim \lambda_W \ll a_0$. \\
Spin & Unknot chirality & Two orientations of the unknot twist: $\pm 1/2$. \\
Exchange integral $K_{12}$ & Phase-separation geometry & For same orbital: $K_{12} = J_{12}$ (mathematics identical). \\
Correlation energy & LC phase-jitter & \textbf{New AVE content.}  $J_{s^2} = \tfrac{1}{2}(1+p_c)$ from torus-knot geometry. \\
\bottomrule
\end{tabular}
\end{table}


\begin{table}[H]
\centering
\caption{Quantum Mechanics $\leftrightarrow$ AVE Translation.  Rows 1--8: same mathematics, different ontology.  Row~9: genuinely new AVE prediction.}
\label{tab:trans_qm}
\small
\renewcommand{\arraystretch}{1.3}
\begin{tabular}{@{}p{3.2cm} p{3.5cm} p{7cm}@{}}
\toprule
\textbf{Standard QM} & \textbf{AVE Equivalent} & \textbf{Relationship} \\
\midrule
Wavefunction $\psi(r)$ & Acoustic cavity mode & Same equation (Helmholtz). AVE: mode of vacuum LC mesh. \\
$|\psi|^2$ probability & Time-averaged trajectory & Density of point-defect sweeping its standing-wave mode. \\
Coulomb potential $e^2/(4\pi\varepsilon_0 r)$ & $\alpha\hbar c / r$ (Axiom 2) & Algebraically identical. $K = \alpha\hbar c$ is the impedance coupling. \\
Bohr radius $a_0$ & $\ell_{node}/\alpha$ & $a_0 = \hbar/(m_e c \alpha)$. The atom is $1/\alpha \approx 137$ lattice pitches. \\
Hartree energy $E_H$ & $m_e(\alpha c)^2$ & Same formula. $E_H = K/a_0$ (coupling per cavity radius). \\
Electron & $0_1$ topological unknot & Irreducible closed flux tube. Internal radius $\sim \lambda_W \ll a_0$. \\
Spin & Unknot chirality & Two orientations of the unknot twist: $\pm 1/2$. \\
Exchange integral $K_{12}$ & Phase-separation geometry & For same orbital: $K_{12} = J_{12}$ (mathematics identical). \\
Correlation energy & LC phase-jitter & \textbf{New AVE content.}  $J_{s^2} = \tfrac{1}{2}(1+p_c)$ from torus-knot geometry. \\
\bottomrule
\end{tabular}
\end{table}


\begin{table}[H]
\centering
\caption{Quantum Mechanics $\leftrightarrow$ AVE Translation.  Rows 1--8: same mathematics, different ontology.  Row~9: genuinely new AVE prediction.}
\label{tab:trans_qm}
\small
\renewcommand{\arraystretch}{1.3}
\begin{tabular}{@{}p{3.2cm} p{3.5cm} p{7cm}@{}}
\toprule
\textbf{Standard QM} & \textbf{AVE Equivalent} & \textbf{Relationship} \\
\midrule
Wavefunction $\psi(r)$ & Acoustic cavity mode & Same equation (Helmholtz). AVE: mode of vacuum LC mesh. \\
$|\psi|^2$ probability & Time-averaged trajectory & Density of point-defect sweeping its standing-wave mode. \\
Coulomb potential $e^2/(4\pi\varepsilon_0 r)$ & $\alpha\hbar c / r$ (Axiom 2) & Algebraically identical. $K = \alpha\hbar c$ is the impedance coupling. \\
Bohr radius $a_0$ & $\ell_{node}/\alpha$ & $a_0 = \hbar/(m_e c \alpha)$. The atom is $1/\alpha \approx 137$ lattice pitches. \\
Hartree energy $E_H$ & $m_e(\alpha c)^2$ & Same formula. $E_H = K/a_0$ (coupling per cavity radius). \\
Electron & $0_1$ topological unknot & Irreducible closed flux tube. Internal radius $\sim \lambda_W \ll a_0$. \\
Spin & Unknot chirality & Two orientations of the unknot twist: $\pm 1/2$. \\
Exchange integral $K_{12}$ & Phase-separation geometry & For same orbital: $K_{12} = J_{12}$ (mathematics identical). \\
Correlation energy & LC phase-jitter & \textbf{New AVE content.}  $J_{s^2} = \tfrac{1}{2}(1+p_c)$ from torus-knot geometry. \\
\bottomrule
\end{tabular}
\end{table}


\subsection{Particle Physics / Standard Model}
% ═══════════════════════════════════════════════════════════════
% DOMAIN TRANSLATION TABLE: Particle Physics / Standard Model ↔ AVE
% Canonical location: manuscript/common/translation_particle_physics.tex
% ═══════════════════════════════════════════════════════════════
%
% Usage: % ═══════════════════════════════════════════════════════════════
% DOMAIN TRANSLATION TABLE: Particle Physics / Standard Model ↔ AVE
% Canonical location: manuscript/common/translation_particle_physics.tex
% ═══════════════════════════════════════════════════════════════
%
% Usage: % ═══════════════════════════════════════════════════════════════
% DOMAIN TRANSLATION TABLE: Particle Physics / Standard Model ↔ AVE
% Canonical location: manuscript/common/translation_particle_physics.tex
% ═══════════════════════════════════════════════════════════════
%
% Usage: \input{../../common/translation_particle_physics.tex}

\begin{table}[H]
\centering
\caption{Standard Model $\leftrightarrow$ AVE Translation.  Particles are topological defects; forces are impedance gradients.}
\label{tab:trans_particle}
\small
\renewcommand{\arraystretch}{1.3}
\begin{tabular}{@{}p{3.8cm} p{4.0cm} p{5.8cm}@{}}
\toprule
\textbf{Standard Model} & \textbf{AVE Equivalent} & \textbf{Relationship} \\
\midrule
Electron & $0_1$ unknot (translation sector) & $m_e = $ topological self-impedance. Axiom~1. \\
Muon, tau & $0_1$ unknots (rotation, curvature) & Three Cosserat sectors $\to$ three generations. \\
Quark & Fractional vortex ($\theta = 2\pi/3$) & Witten effect: $q = n + \theta e/(2\pi)$. \\
Gluon & Confining flux tube & $F_{conf} = 3(m_p/m_e)\alpha^{-1}T_{EM} \approx 1$~GeV/fm. \\
Photon & Transverse shear wave & Linear acoustic mode of the LC mesh. \\
$W^\pm, Z^0$ bosons & Torsional gauge modes & $m_W/m_Z = \sqrt{7}/3$ from $\nu_{vac} = 2/7$. \\
Higgs field & Dielectric saturation (Axiom~4) & $v = 1/\sqrt{\sqrt{2}\,G_F} \approx 248.8$~GeV; spontaneous yield. \\
$\sin^2\theta_W$ & Acoustic mode ratio & $1 - m_W^2/m_Z^2 = 2/9$. \\
$\alpha_s$ (strong coupling) & Geometric projection & $\alpha_s = \alpha^{3/7}$; 3 spatial of 7 total DOF. \\
Proton mass & Faddeev-Skyrme eigenvalue & $m_p \approx 1836\,m_e$ from $\kappa_{FS}, p_c$. \\
Neutrino mass & Torsional-to-translational ratio & $m_\nu \approx \alpha(m_e/M_W)m_e \sim 0.024$~eV. \\
CKM matrix & Dielectric mixing cascade & $V_{us} = \sin\theta_C = 2/9$ from $\sin^2\theta_W$. \\
Empirical Mass Defect (8~MeV) & $\alpha M_p c^2$ Yield Limit & Binding ceiling is purely the Axiom 4 rupture limit ($\sim 6.84$~MeV). \\
Isotopic Curve of Stability & Geometric Coulomb Penalty & Required neutron packing scales via topological limit $N \approx Z + 1.2\alpha Z^2$. \\
\bottomrule
\end{tabular}
\end{table}


\begin{table}[H]
\centering
\caption{Standard Model $\leftrightarrow$ AVE Translation.  Particles are topological defects; forces are impedance gradients.}
\label{tab:trans_particle}
\small
\renewcommand{\arraystretch}{1.3}
\begin{tabular}{@{}p{3.8cm} p{4.0cm} p{5.8cm}@{}}
\toprule
\textbf{Standard Model} & \textbf{AVE Equivalent} & \textbf{Relationship} \\
\midrule
Electron & $0_1$ unknot (translation sector) & $m_e = $ topological self-impedance. Axiom~1. \\
Muon, tau & $0_1$ unknots (rotation, curvature) & Three Cosserat sectors $\to$ three generations. \\
Quark & Fractional vortex ($\theta = 2\pi/3$) & Witten effect: $q = n + \theta e/(2\pi)$. \\
Gluon & Confining flux tube & $F_{conf} = 3(m_p/m_e)\alpha^{-1}T_{EM} \approx 1$~GeV/fm. \\
Photon & Transverse shear wave & Linear acoustic mode of the LC mesh. \\
$W^\pm, Z^0$ bosons & Torsional gauge modes & $m_W/m_Z = \sqrt{7}/3$ from $\nu_{vac} = 2/7$. \\
Higgs field & Dielectric saturation (Axiom~4) & $v = 1/\sqrt{\sqrt{2}\,G_F} \approx 248.8$~GeV; spontaneous yield. \\
$\sin^2\theta_W$ & Acoustic mode ratio & $1 - m_W^2/m_Z^2 = 2/9$. \\
$\alpha_s$ (strong coupling) & Geometric projection & $\alpha_s = \alpha^{3/7}$; 3 spatial of 7 total DOF. \\
Proton mass & Faddeev-Skyrme eigenvalue & $m_p \approx 1836\,m_e$ from $\kappa_{FS}, p_c$. \\
Neutrino mass & Torsional-to-translational ratio & $m_\nu \approx \alpha(m_e/M_W)m_e \sim 0.024$~eV. \\
CKM matrix & Dielectric mixing cascade & $V_{us} = \sin\theta_C = 2/9$ from $\sin^2\theta_W$. \\
Empirical Mass Defect (8~MeV) & $\alpha M_p c^2$ Yield Limit & Binding ceiling is purely the Axiom 4 rupture limit ($\sim 6.84$~MeV). \\
Isotopic Curve of Stability & Geometric Coulomb Penalty & Required neutron packing scales via topological limit $N \approx Z + 1.2\alpha Z^2$. \\
\bottomrule
\end{tabular}
\end{table}


\begin{table}[H]
\centering
\caption{Standard Model $\leftrightarrow$ AVE Translation.  Particles are topological defects; forces are impedance gradients.}
\label{tab:trans_particle}
\small
\renewcommand{\arraystretch}{1.3}
\begin{tabular}{@{}p{3.8cm} p{4.0cm} p{5.8cm}@{}}
\toprule
\textbf{Standard Model} & \textbf{AVE Equivalent} & \textbf{Relationship} \\
\midrule
Electron & $0_1$ unknot (translation sector) & $m_e = $ topological self-impedance. Axiom~1. \\
Muon, tau & $0_1$ unknots (rotation, curvature) & Three Cosserat sectors $\to$ three generations. \\
Quark & Fractional vortex ($\theta = 2\pi/3$) & Witten effect: $q = n + \theta e/(2\pi)$. \\
Gluon & Confining flux tube & $F_{conf} = 3(m_p/m_e)\alpha^{-1}T_{EM} \approx 1$~GeV/fm. \\
Photon & Transverse shear wave & Linear acoustic mode of the LC mesh. \\
$W^\pm, Z^0$ bosons & Torsional gauge modes & $m_W/m_Z = \sqrt{7}/3$ from $\nu_{vac} = 2/7$. \\
Higgs field & Dielectric saturation (Axiom~4) & $v = 1/\sqrt{\sqrt{2}\,G_F} \approx 248.8$~GeV; spontaneous yield. \\
$\sin^2\theta_W$ & Acoustic mode ratio & $1 - m_W^2/m_Z^2 = 2/9$. \\
$\alpha_s$ (strong coupling) & Geometric projection & $\alpha_s = \alpha^{3/7}$; 3 spatial of 7 total DOF. \\
Proton mass & Faddeev-Skyrme eigenvalue & $m_p \approx 1836\,m_e$ from $\kappa_{FS}, p_c$. \\
Neutrino mass & Torsional-to-translational ratio & $m_\nu \approx \alpha(m_e/M_W)m_e \sim 0.024$~eV. \\
CKM matrix & Dielectric mixing cascade & $V_{us} = \sin\theta_C = 2/9$ from $\sin^2\theta_W$. \\
Empirical Mass Defect (8~MeV) & $\alpha M_p c^2$ Yield Limit & Binding ceiling is purely the Axiom 4 rupture limit ($\sim 6.84$~MeV). \\
Isotopic Curve of Stability & Geometric Coulomb Penalty & Required neutron packing scales via topological limit $N \approx Z + 1.2\alpha Z^2$. \\
\bottomrule
\end{tabular}
\end{table}


\subsection{General Relativity / Gravity}
% ═══════════════════════════════════════════════════════════════
% DOMAIN TRANSLATION TABLE: General Relativity / Gravity ↔ AVE
% Canonical location: manuscript/common/translation_gravity.tex
% ═══════════════════════════════════════════════════════════════
%
% Usage: % ═══════════════════════════════════════════════════════════════
% DOMAIN TRANSLATION TABLE: General Relativity / Gravity ↔ AVE
% Canonical location: manuscript/common/translation_gravity.tex
% ═══════════════════════════════════════════════════════════════
%
% Usage: % ═══════════════════════════════════════════════════════════════
% DOMAIN TRANSLATION TABLE: General Relativity / Gravity ↔ AVE
% Canonical location: manuscript/common/translation_gravity.tex
% ═══════════════════════════════════════════════════════════════
%
% Usage: \input{../../common/translation_gravity.tex}

\begin{table}[H]
\centering
\caption{General Relativity $\leftrightarrow$ AVE Translation.  Gravity is macroscopic dielectric refraction; no curved manifold is required.}
\label{tab:trans_gravity}
\small
\renewcommand{\arraystretch}{1.3}
\begin{tabular}{@{}p{3.8cm} p{4.0cm} p{5.8cm}@{}}
\toprule
\textbf{General Relativity} & \textbf{AVE Equivalent} & \textbf{Relationship} \\
\midrule
Metric tensor $g_{\mu\nu}$ & Variable $\varepsilon_{eff}, \mu_{eff}$ & $n(r) = 1 + 2GM/(rc^2)$; symmetric impedance gradient. \\
Spacetime curvature & Dielectric refraction & Geodesic $\equiv$ Fermat path through refractive medium. \\
Stress-energy $T_{\mu\nu}$ & LC energy density $U$ & $U = \tfrac{1}{2}\varepsilon_0 |E|^2 + \tfrac{1}{2}\mu_0 |H|^2$. \\
Schwarzschild radius $r_s$ & Saturation boundary ($S \to 0$) & $Z \to 0, \; \Gamma = -1$: total impedance collapse. \\
Event horizon & Catastrophic impedance boundary & Axiom~4: dielectric yield $\to$ total confinement. \\
Gravitational waves & Transverse inductive shear waves & Low-frequency shear modes in the LC condensate. \\
Frame dragging (Kerr) & Acoustic vortex flow $\vec{v}_\phi$ & Gordon metric: $\omega(r) = 2Mar/(r^2+a^2)^2$. \\
Gravitational lensing & Optical refraction ($n > 1$) & $n(r) = (1+r_s/2r)^3/(1-r_s/2r)$. \\
$G$ (Newton) & TT strain projection & $G = \alpha^2 c^3 \ell_{node}/(28\pi m_e)$ (Axiom~2). \\
``Dark matter'' halo & Regime~I vacuum drag & $v_{flat} = (GM a_0)^{1/4}$; $a_0 = cH_\infty/(2\pi)$. \\
\bottomrule
\end{tabular}
\end{table}


\begin{table}[H]
\centering
\caption{General Relativity $\leftrightarrow$ AVE Translation.  Gravity is macroscopic dielectric refraction; no curved manifold is required.}
\label{tab:trans_gravity}
\small
\renewcommand{\arraystretch}{1.3}
\begin{tabular}{@{}p{3.8cm} p{4.0cm} p{5.8cm}@{}}
\toprule
\textbf{General Relativity} & \textbf{AVE Equivalent} & \textbf{Relationship} \\
\midrule
Metric tensor $g_{\mu\nu}$ & Variable $\varepsilon_{eff}, \mu_{eff}$ & $n(r) = 1 + 2GM/(rc^2)$; symmetric impedance gradient. \\
Spacetime curvature & Dielectric refraction & Geodesic $\equiv$ Fermat path through refractive medium. \\
Stress-energy $T_{\mu\nu}$ & LC energy density $U$ & $U = \tfrac{1}{2}\varepsilon_0 |E|^2 + \tfrac{1}{2}\mu_0 |H|^2$. \\
Schwarzschild radius $r_s$ & Saturation boundary ($S \to 0$) & $Z \to 0, \; \Gamma = -1$: total impedance collapse. \\
Event horizon & Catastrophic impedance boundary & Axiom~4: dielectric yield $\to$ total confinement. \\
Gravitational waves & Transverse inductive shear waves & Low-frequency shear modes in the LC condensate. \\
Frame dragging (Kerr) & Acoustic vortex flow $\vec{v}_\phi$ & Gordon metric: $\omega(r) = 2Mar/(r^2+a^2)^2$. \\
Gravitational lensing & Optical refraction ($n > 1$) & $n(r) = (1+r_s/2r)^3/(1-r_s/2r)$. \\
$G$ (Newton) & TT strain projection & $G = \alpha^2 c^3 \ell_{node}/(28\pi m_e)$ (Axiom~2). \\
``Dark matter'' halo & Regime~I vacuum drag & $v_{flat} = (GM a_0)^{1/4}$; $a_0 = cH_\infty/(2\pi)$. \\
\bottomrule
\end{tabular}
\end{table}


\begin{table}[H]
\centering
\caption{General Relativity $\leftrightarrow$ AVE Translation.  Gravity is macroscopic dielectric refraction; no curved manifold is required.}
\label{tab:trans_gravity}
\small
\renewcommand{\arraystretch}{1.3}
\begin{tabular}{@{}p{3.8cm} p{4.0cm} p{5.8cm}@{}}
\toprule
\textbf{General Relativity} & \textbf{AVE Equivalent} & \textbf{Relationship} \\
\midrule
Metric tensor $g_{\mu\nu}$ & Variable $\varepsilon_{eff}, \mu_{eff}$ & $n(r) = 1 + 2GM/(rc^2)$; symmetric impedance gradient. \\
Spacetime curvature & Dielectric refraction & Geodesic $\equiv$ Fermat path through refractive medium. \\
Stress-energy $T_{\mu\nu}$ & LC energy density $U$ & $U = \tfrac{1}{2}\varepsilon_0 |E|^2 + \tfrac{1}{2}\mu_0 |H|^2$. \\
Schwarzschild radius $r_s$ & Saturation boundary ($S \to 0$) & $Z \to 0, \; \Gamma = -1$: total impedance collapse. \\
Event horizon & Catastrophic impedance boundary & Axiom~4: dielectric yield $\to$ total confinement. \\
Gravitational waves & Transverse inductive shear waves & Low-frequency shear modes in the LC condensate. \\
Frame dragging (Kerr) & Acoustic vortex flow $\vec{v}_\phi$ & Gordon metric: $\omega(r) = 2Mar/(r^2+a^2)^2$. \\
Gravitational lensing & Optical refraction ($n > 1$) & $n(r) = (1+r_s/2r)^3/(1-r_s/2r)$. \\
$G$ (Newton) & TT strain projection & $G = \alpha^2 c^3 \ell_{node}/(28\pi m_e)$ (Axiom~2). \\
``Dark matter'' halo & Regime~I vacuum drag & $v_{flat} = (GM a_0)^{1/4}$; $a_0 = cH_\infty/(2\pi)$. \\
\bottomrule
\end{tabular}
\end{table}


\subsection{Cosmology}
% ═══════════════════════════════════════════════════════════════
% DOMAIN TRANSLATION TABLE: Cosmology ↔ AVE
% Canonical location: manuscript/common/translation_cosmology.tex
% ═══════════════════════════════════════════════════════════════
%
% Usage: % ═══════════════════════════════════════════════════════════════
% DOMAIN TRANSLATION TABLE: Cosmology ↔ AVE
% Canonical location: manuscript/common/translation_cosmology.tex
% ═══════════════════════════════════════════════════════════════
%
% Usage: % ═══════════════════════════════════════════════════════════════
% DOMAIN TRANSLATION TABLE: Cosmology ↔ AVE
% Canonical location: manuscript/common/translation_cosmology.tex
% ═══════════════════════════════════════════════════════════════
%
% Usage: \input{../../common/translation_cosmology.tex}

\begin{table}[H]
\centering
\caption{Cosmology $\leftrightarrow$ AVE Translation. Expansion, dark energy, and dark matter emerge from lattice creep and impedance gradients.}
\label{tab:trans_cosmology}
\small
\renewcommand{\arraystretch}{1.3}
\begin{tabular}{@{}p{3.8cm} p{4.0cm} p{5.8cm}@{}}
\toprule
\textbf{Cosmology} & \textbf{AVE Equivalent} & \textbf{Relationship} \\
\midrule
Hubble expansion $H_0$ & Lattice creep rate $H_\infty$ & $H_\infty = 28\pi m_e^3 cG/(\hbar^2\alpha^2) \approx 69.3$~km/s/Mpc. \\
Dark energy ($\Lambda$) & Latent strain energy & $w_{vac} < -1$: stable phantom from lattice pre-stress. \\
Dark matter & Regime~I vacuum drag & Unsaturated lattice mutual inductance below $a_0$. \\
MOND acceleration $a_0$ & Regime boundary & $a_0 = cH_\infty/(2\pi) \approx 1.07 \times 10^{-10}$~m/s$^2$. \\
Flat rotation curves & Asymptotic visco-kinematic floor & $v_{flat} = (GM_{bar}\,a_0)^{1/4}$ from hoop stress. \\
CMB temperature & Thermal lattice noise $T_{CMB}$ & Residual acoustic background of the LC mesh. \\
Horizon scale $R_H$ & $\alpha^2/(28\pi\alpha_G)$ lattice pitches & $R_H \approx 14.1$~Gly. \\
Inflation & Not required & Lattice isotropy is built in (Axiom~1). \\
Baryon asymmetry & Topological chirality bias & $\mathcal{M}_A$ chiral LC lattice breaks CP natively. \\
\bottomrule
\end{tabular}
\end{table}


\begin{table}[H]
\centering
\caption{Cosmology $\leftrightarrow$ AVE Translation. Expansion, dark energy, and dark matter emerge from lattice creep and impedance gradients.}
\label{tab:trans_cosmology}
\small
\renewcommand{\arraystretch}{1.3}
\begin{tabular}{@{}p{3.8cm} p{4.0cm} p{5.8cm}@{}}
\toprule
\textbf{Cosmology} & \textbf{AVE Equivalent} & \textbf{Relationship} \\
\midrule
Hubble expansion $H_0$ & Lattice creep rate $H_\infty$ & $H_\infty = 28\pi m_e^3 cG/(\hbar^2\alpha^2) \approx 69.3$~km/s/Mpc. \\
Dark energy ($\Lambda$) & Latent strain energy & $w_{vac} < -1$: stable phantom from lattice pre-stress. \\
Dark matter & Regime~I vacuum drag & Unsaturated lattice mutual inductance below $a_0$. \\
MOND acceleration $a_0$ & Regime boundary & $a_0 = cH_\infty/(2\pi) \approx 1.07 \times 10^{-10}$~m/s$^2$. \\
Flat rotation curves & Asymptotic visco-kinematic floor & $v_{flat} = (GM_{bar}\,a_0)^{1/4}$ from hoop stress. \\
CMB temperature & Thermal lattice noise $T_{CMB}$ & Residual acoustic background of the LC mesh. \\
Horizon scale $R_H$ & $\alpha^2/(28\pi\alpha_G)$ lattice pitches & $R_H \approx 14.1$~Gly. \\
Inflation & Not required & Lattice isotropy is built in (Axiom~1). \\
Baryon asymmetry & Topological chirality bias & $\mathcal{M}_A$ chiral LC lattice breaks CP natively. \\
\bottomrule
\end{tabular}
\end{table}


\begin{table}[H]
\centering
\caption{Cosmology $\leftrightarrow$ AVE Translation. Expansion, dark energy, and dark matter emerge from lattice creep and impedance gradients.}
\label{tab:trans_cosmology}
\small
\renewcommand{\arraystretch}{1.3}
\begin{tabular}{@{}p{3.8cm} p{4.0cm} p{5.8cm}@{}}
\toprule
\textbf{Cosmology} & \textbf{AVE Equivalent} & \textbf{Relationship} \\
\midrule
Hubble expansion $H_0$ & Lattice creep rate $H_\infty$ & $H_\infty = 28\pi m_e^3 cG/(\hbar^2\alpha^2) \approx 69.3$~km/s/Mpc. \\
Dark energy ($\Lambda$) & Latent strain energy & $w_{vac} < -1$: stable phantom from lattice pre-stress. \\
Dark matter & Regime~I vacuum drag & Unsaturated lattice mutual inductance below $a_0$. \\
MOND acceleration $a_0$ & Regime boundary & $a_0 = cH_\infty/(2\pi) \approx 1.07 \times 10^{-10}$~m/s$^2$. \\
Flat rotation curves & Asymptotic visco-kinematic floor & $v_{flat} = (GM_{bar}\,a_0)^{1/4}$ from hoop stress. \\
CMB temperature & Thermal lattice noise $T_{CMB}$ & Residual acoustic background of the LC mesh. \\
Horizon scale $R_H$ & $\alpha^2/(28\pi\alpha_G)$ lattice pitches & $R_H \approx 14.1$~Gly. \\
Inflation & Not required & Lattice isotropy is built in (Axiom~1). \\
Baryon asymmetry & Topological chirality bias & $\mathcal{M}_A$ chiral LC lattice breaks CP natively. \\
\bottomrule
\end{tabular}
\end{table}


\subsection{Condensed Matter / Superconductivity}
% ═══════════════════════════════════════════════════════════════
% DOMAIN TRANSLATION TABLE: Condensed Matter ↔ AVE
% Canonical location: manuscript/common/translation_condensed_matter.tex
% ═══════════════════════════════════════════════════════════════
%
% Usage: % ═══════════════════════════════════════════════════════════════
% DOMAIN TRANSLATION TABLE: Condensed Matter ↔ AVE
% Canonical location: manuscript/common/translation_condensed_matter.tex
% ═══════════════════════════════════════════════════════════════
%
% Usage: % ═══════════════════════════════════════════════════════════════
% DOMAIN TRANSLATION TABLE: Condensed Matter ↔ AVE
% Canonical location: manuscript/common/translation_condensed_matter.tex
% ═══════════════════════════════════════════════════════════════
%
% Usage: \input{../../common/translation_condensed_matter.tex}

\begin{table}[H]
\centering
\caption{Condensed Matter $\leftrightarrow$ AVE Translation.  Superconductivity is classical Kuramoto phase-locking of topological inductors.}
\label{tab:trans_condensed}
\small
\renewcommand{\arraystretch}{1.3}
\begin{tabular}{@{}p{3.8cm} p{4.0cm} p{5.8cm}@{}}
\toprule
\textbf{Condensed Matter} & \textbf{AVE Equivalent} & \textbf{Relationship} \\
\midrule
Cooper pair & Phase-locked unknot pair & Kuramoto synchronisation at $T < T_c$. \\
BCS condensate & Macroscopic gear train & All $0_1$ unknots locked: $\dot{\theta}_i = \Omega_{macro}$. \\
Superconductivity ($R=0$) & Zero relative induction & $\Delta(dB/dt)_{rel} = 0$ when phase-locked. \\
Critical temperature $T_c$ & Coupling threshold & Thermal jitter $\xi(T) < K$ (Kuramoto coupling). \\
Meissner effect & Boundary torque rejection & Interlocked gyroscope inertia expels external $B$. \\
London penetration $\lambda_L$ & Inertial decay length & $\omega(x) = \omega_0 e^{-x/\lambda_\text{inertial}}$. \\
$B_c(T)$ critical field & Axiom~4 saturation & $B_c = B_{c0}\sqrt{1-(T/T_c)^2} \equiv S(T/T_c)$. \\
Phonon mediation & Lattice acoustic coupling & Same transverse shear modes; no virtual exchange. \\
Type~I / Type~II & Regime~I / I--II boundary & Saturation depth vs.\ coherence length. \\
Band gap & Topological mass gap & $\Delta = \kappa_{FS} \cdot \ell_{node}$ coupling energy. \\
Fluid PVT Phase State & LC Impedance Partition Fraction & Macroscopic anomalies (e.g., $4^\circ$C water) driven by intact tetrahedral H-bond vs thermal shatter. \\
\bottomrule
\end{tabular}
\end{table}


\begin{table}[H]
\centering
\caption{Condensed Matter $\leftrightarrow$ AVE Translation.  Superconductivity is classical Kuramoto phase-locking of topological inductors.}
\label{tab:trans_condensed}
\small
\renewcommand{\arraystretch}{1.3}
\begin{tabular}{@{}p{3.8cm} p{4.0cm} p{5.8cm}@{}}
\toprule
\textbf{Condensed Matter} & \textbf{AVE Equivalent} & \textbf{Relationship} \\
\midrule
Cooper pair & Phase-locked unknot pair & Kuramoto synchronisation at $T < T_c$. \\
BCS condensate & Macroscopic gear train & All $0_1$ unknots locked: $\dot{\theta}_i = \Omega_{macro}$. \\
Superconductivity ($R=0$) & Zero relative induction & $\Delta(dB/dt)_{rel} = 0$ when phase-locked. \\
Critical temperature $T_c$ & Coupling threshold & Thermal jitter $\xi(T) < K$ (Kuramoto coupling). \\
Meissner effect & Boundary torque rejection & Interlocked gyroscope inertia expels external $B$. \\
London penetration $\lambda_L$ & Inertial decay length & $\omega(x) = \omega_0 e^{-x/\lambda_\text{inertial}}$. \\
$B_c(T)$ critical field & Axiom~4 saturation & $B_c = B_{c0}\sqrt{1-(T/T_c)^2} \equiv S(T/T_c)$. \\
Phonon mediation & Lattice acoustic coupling & Same transverse shear modes; no virtual exchange. \\
Type~I / Type~II & Regime~I / I--II boundary & Saturation depth vs.\ coherence length. \\
Band gap & Topological mass gap & $\Delta = \kappa_{FS} \cdot \ell_{node}$ coupling energy. \\
Fluid PVT Phase State & LC Impedance Partition Fraction & Macroscopic anomalies (e.g., $4^\circ$C water) driven by intact tetrahedral H-bond vs thermal shatter. \\
\bottomrule
\end{tabular}
\end{table}


\begin{table}[H]
\centering
\caption{Condensed Matter $\leftrightarrow$ AVE Translation.  Superconductivity is classical Kuramoto phase-locking of topological inductors.}
\label{tab:trans_condensed}
\small
\renewcommand{\arraystretch}{1.3}
\begin{tabular}{@{}p{3.8cm} p{4.0cm} p{5.8cm}@{}}
\toprule
\textbf{Condensed Matter} & \textbf{AVE Equivalent} & \textbf{Relationship} \\
\midrule
Cooper pair & Phase-locked unknot pair & Kuramoto synchronisation at $T < T_c$. \\
BCS condensate & Macroscopic gear train & All $0_1$ unknots locked: $\dot{\theta}_i = \Omega_{macro}$. \\
Superconductivity ($R=0$) & Zero relative induction & $\Delta(dB/dt)_{rel} = 0$ when phase-locked. \\
Critical temperature $T_c$ & Coupling threshold & Thermal jitter $\xi(T) < K$ (Kuramoto coupling). \\
Meissner effect & Boundary torque rejection & Interlocked gyroscope inertia expels external $B$. \\
London penetration $\lambda_L$ & Inertial decay length & $\omega(x) = \omega_0 e^{-x/\lambda_\text{inertial}}$. \\
$B_c(T)$ critical field & Axiom~4 saturation & $B_c = B_{c0}\sqrt{1-(T/T_c)^2} \equiv S(T/T_c)$. \\
Phonon mediation & Lattice acoustic coupling & Same transverse shear modes; no virtual exchange. \\
Type~I / Type~II & Regime~I / I--II boundary & Saturation depth vs.\ coherence length. \\
Band gap & Topological mass gap & $\Delta = \kappa_{FS} \cdot \ell_{node}$ coupling energy. \\
Fluid PVT Phase State & LC Impedance Partition Fraction & Macroscopic anomalies (e.g., $4^\circ$C water) driven by intact tetrahedral H-bond vs thermal shatter. \\
\bottomrule
\end{tabular}
\end{table}


\subsection{Protein Folding (Biology $\leftrightarrow$ EE)}
% ═══════════════════════════════════════════════════════════════
% DOMAIN TRANSLATION TABLE: Protein Folding (Biology ↔ EE ↔ AVE)
% Source: Extracted from Book 6 Ch.2
% Canonical location: manuscript/common/translation_protein.tex
% ═══════════════════════════════════════════════════════════════
%
% Usage: % ═══════════════════════════════════════════════════════════════
% DOMAIN TRANSLATION TABLE: Protein Folding (Biology ↔ EE ↔ AVE)
% Source: Extracted from Book 6 Ch.2
% Canonical location: manuscript/common/translation_protein.tex
% ═══════════════════════════════════════════════════════════════
%
% Usage: % ═══════════════════════════════════════════════════════════════
% DOMAIN TRANSLATION TABLE: Protein Folding (Biology ↔ EE ↔ AVE)
% Source: Extracted from Book 6 Ch.2
% Canonical location: manuscript/common/translation_protein.tex
% ═══════════════════════════════════════════════════════════════
%
% Usage: \input{../../common/translation_protein.tex}

\begin{table}[H]
\centering
\caption{Protein folding terminology: structural biology $\leftrightarrow$ transmission line / circuit theory $\leftrightarrow$ AVE derivation.}
\label{tab:trans_protein}
\small
\renewcommand{\arraystretch}{1.2}
\begin{tabular}{@{}p{4.5cm} p{4.5cm} p{4.5cm}@{}}
\toprule
\textbf{Biology} & \textbf{Electrical Engineering} & \textbf{AVE Derivation} \\
\midrule
Amino acid backbone & Cascaded transmission line & $L$-$C$ ladder from bond stiffness \\
Sidechain R-group & Shunt stub impedance & $Z_R(\omega)$ at backbone junction \\
Helix propensity $P_\alpha$ & Impedance match ($Z_R \gg Z_0$) & Low $Z_{topo}$: stub invisible \\
Sheet / coil tendency & Impedance mismatch ($Z_R \sim Z_0$) & High $Z_{topo}$: destructive interference \\
Ramachandran map & Smith chart & Allowed $\Gamma$ trajectories \\
Steric clash & Short-circuit reflection ($\Gamma = -1$) & $d < r_A + r_B$ (Pauli exclusion) \\
H-bond ($i \to i{+}4$) & Series resonant coupling & $L$-$C$ energy transfer at $\omega_0$ \\
Inter-strand H-bond & Coupled microstrip lines & Mutual $L$/$C$ at $d_\beta = 4.7$~\AA \\
Hydrophobic effect & Impedance mismatch with termination & $h_i \cdot h_j$ coupling (water $\varepsilon_r \approx 80$) \\
Native fold & Matched load (zero reflection) & Global $U_{\text{total}}$ minimum \\
Levinthal's paradox & Why doesn't the line ring forever? & Deterministic $Z$-driven gradient \\
Disulfide bond & Topological short-circuit & Non-local $Z$-match creates loop \\
Salt bridge (Glu--Lys) & Conjugate impedance match & $\operatorname{Re}(Z_i \cdot Z_j^*) > 0$ \\
Solvent (water) & Chassis ground (parasitic shunt) & $Y_{\text{solv}} = \text{exposure}_i / Z_{\text{H}_2\text{O}}$ \\
L-amino acid chirality & Non-reciprocal waveguide & $\beta_{\text{eff}} = \beta_0 - \delta_\chi \tanh(\chi/\chi_0)$ \\
$\alpha$-helix & Helical slow-wave structure & Right-handed coiled delay line \\
$\beta$-sheet & Coupled stripline array & Antiparallel microstrip pair \\
\bottomrule
\end{tabular}
\end{table}


\begin{table}[H]
\centering
\caption{Protein folding terminology: structural biology $\leftrightarrow$ transmission line / circuit theory $\leftrightarrow$ AVE derivation.}
\label{tab:trans_protein}
\small
\renewcommand{\arraystretch}{1.2}
\begin{tabular}{@{}p{4.5cm} p{4.5cm} p{4.5cm}@{}}
\toprule
\textbf{Biology} & \textbf{Electrical Engineering} & \textbf{AVE Derivation} \\
\midrule
Amino acid backbone & Cascaded transmission line & $L$-$C$ ladder from bond stiffness \\
Sidechain R-group & Shunt stub impedance & $Z_R(\omega)$ at backbone junction \\
Helix propensity $P_\alpha$ & Impedance match ($Z_R \gg Z_0$) & Low $Z_{topo}$: stub invisible \\
Sheet / coil tendency & Impedance mismatch ($Z_R \sim Z_0$) & High $Z_{topo}$: destructive interference \\
Ramachandran map & Smith chart & Allowed $\Gamma$ trajectories \\
Steric clash & Short-circuit reflection ($\Gamma = -1$) & $d < r_A + r_B$ (Pauli exclusion) \\
H-bond ($i \to i{+}4$) & Series resonant coupling & $L$-$C$ energy transfer at $\omega_0$ \\
Inter-strand H-bond & Coupled microstrip lines & Mutual $L$/$C$ at $d_\beta = 4.7$~\AA \\
Hydrophobic effect & Impedance mismatch with termination & $h_i \cdot h_j$ coupling (water $\varepsilon_r \approx 80$) \\
Native fold & Matched load (zero reflection) & Global $U_{\text{total}}$ minimum \\
Levinthal's paradox & Why doesn't the line ring forever? & Deterministic $Z$-driven gradient \\
Disulfide bond & Topological short-circuit & Non-local $Z$-match creates loop \\
Salt bridge (Glu--Lys) & Conjugate impedance match & $\operatorname{Re}(Z_i \cdot Z_j^*) > 0$ \\
Solvent (water) & Chassis ground (parasitic shunt) & $Y_{\text{solv}} = \text{exposure}_i / Z_{\text{H}_2\text{O}}$ \\
L-amino acid chirality & Non-reciprocal waveguide & $\beta_{\text{eff}} = \beta_0 - \delta_\chi \tanh(\chi/\chi_0)$ \\
$\alpha$-helix & Helical slow-wave structure & Right-handed coiled delay line \\
$\beta$-sheet & Coupled stripline array & Antiparallel microstrip pair \\
\bottomrule
\end{tabular}
\end{table}


\begin{table}[H]
\centering
\caption{Protein folding terminology: structural biology $\leftrightarrow$ transmission line / circuit theory $\leftrightarrow$ AVE derivation.}
\label{tab:trans_protein}
\small
\renewcommand{\arraystretch}{1.2}
\begin{tabular}{@{}p{4.5cm} p{4.5cm} p{4.5cm}@{}}
\toprule
\textbf{Biology} & \textbf{Electrical Engineering} & \textbf{AVE Derivation} \\
\midrule
Amino acid backbone & Cascaded transmission line & $L$-$C$ ladder from bond stiffness \\
Sidechain R-group & Shunt stub impedance & $Z_R(\omega)$ at backbone junction \\
Helix propensity $P_\alpha$ & Impedance match ($Z_R \gg Z_0$) & Low $Z_{topo}$: stub invisible \\
Sheet / coil tendency & Impedance mismatch ($Z_R \sim Z_0$) & High $Z_{topo}$: destructive interference \\
Ramachandran map & Smith chart & Allowed $\Gamma$ trajectories \\
Steric clash & Short-circuit reflection ($\Gamma = -1$) & $d < r_A + r_B$ (Pauli exclusion) \\
H-bond ($i \to i{+}4$) & Series resonant coupling & $L$-$C$ energy transfer at $\omega_0$ \\
Inter-strand H-bond & Coupled microstrip lines & Mutual $L$/$C$ at $d_\beta = 4.7$~\AA \\
Hydrophobic effect & Impedance mismatch with termination & $h_i \cdot h_j$ coupling (water $\varepsilon_r \approx 80$) \\
Native fold & Matched load (zero reflection) & Global $U_{\text{total}}$ minimum \\
Levinthal's paradox & Why doesn't the line ring forever? & Deterministic $Z$-driven gradient \\
Disulfide bond & Topological short-circuit & Non-local $Z$-match creates loop \\
Salt bridge (Glu--Lys) & Conjugate impedance match & $\operatorname{Re}(Z_i \cdot Z_j^*) > 0$ \\
Solvent (water) & Chassis ground (parasitic shunt) & $Y_{\text{solv}} = \text{exposure}_i / Z_{\text{H}_2\text{O}}$ \\
L-amino acid chirality & Non-reciprocal waveguide & $\beta_{\text{eff}} = \beta_0 - \delta_\chi \tanh(\chi/\chi_0)$ \\
$\alpha$-helix & Helical slow-wave structure & Right-handed coiled delay line \\
$\beta$-sheet & Coupled stripline array & Antiparallel microstrip pair \\
\bottomrule
\end{tabular}
\end{table}


\subsection{Protein Solver (Biology $\to$ EE $\to$ AVE Axiom)}
% ═══════════════════════════════════════════════════════════════
% DOMAIN TRANSLATION TABLE: Protein Solver Methods
% Source: Extracted from Book 6 Ch.4
% Canonical location: manuscript/common/translation_protein_solver.tex
% ═══════════════════════════════════════════════════════════════
%
% Usage: % ═══════════════════════════════════════════════════════════════
% DOMAIN TRANSLATION TABLE: Protein Solver Methods
% Source: Extracted from Book 6 Ch.4
% Canonical location: manuscript/common/translation_protein_solver.tex
% ═══════════════════════════════════════════════════════════════
%
% Usage: % ═══════════════════════════════════════════════════════════════
% DOMAIN TRANSLATION TABLE: Protein Solver Methods
% Source: Extracted from Book 6 Ch.4
% Canonical location: manuscript/common/translation_protein_solver.tex
% ═══════════════════════════════════════════════════════════════
%
% Usage: \input{../../common/translation_protein_solver.tex}

\begin{table}[H]
\centering
\caption{Protein solver domain translation.  Every row maps one biological concept through EE to the axiomatic source and its codebase implementation.}
\label{tab:trans_protein_solver}
\scriptsize
\renewcommand{\arraystretch}{1.2}
\begin{tabularx}{\textwidth}{@{}X X X X@{}}
\toprule
\textbf{Biology} & \textbf{EE / RF} & \textbf{AVE Axiom} & \textbf{Script Reference} \\
\midrule
Peptide bond & TL segment & Axiom 1: $Z = \sqrt{\mu/\varepsilon}$ & \texttt{backbone\_segments} \\
Amino acid sidechain & Impedance stub ($Z_\text{topo}$) & Axiom 1: $\sqrt{M/n_e}$ & \texttt{protein\_\allowbreak bond\_\allowbreak constants} \\
H-bond & Mutual inductance ($\kappa$) & Axiom 1: transformer coupling & \texttt{dc\_analysis()} \\
$\beta$-sheet & Backward-wave coupler & Axiom 1: antiparallel TL & \texttt{dc\_analysis()} \\
Hydrophobic core & Conjugate impedance match & Axiom 1: $\operatorname{Re}(Z_iZ_j^*)$ & \texttt{dc\_analysis()} \\
Salt bridge & Reactive LC resonance & Axiom 1: $+jX \times -jX$ & \texttt{dc\_analysis()} \\
Disulfide bond & Near-infinite admittance & Axiom 4: $Z \to \infty$ wall & \texttt{dc\_analysis()} \\
Solvent exposure & Shunt to ground via $Z_\text{water}$ & Axiom 2: Debye relaxation & \texttt{dc\_analysis()} \\
Steric clash & Pauli exclusion ($r < r_\text{steric}$) & Axiom 4: saturation wall & \texttt{dc\_analysis()} \\
Protein compaction & Standing wave pattern & Axiom 4: $\eta_\text{eq} = P_C(1-\nu)$ & \texttt{ac\_analysis()} \\
\midrule
\multicolumn{4}{@{}l}{\textit{Solver Methods:}} \\
\midrule
Native fold & Eigenstate ($\lambda_{\min}(S^\dagger S) = 0$) & 5-step eigenvalue method & \texttt{\_eigenvalue\_target()} \\
Folding kinetics & SPICE transient ring-down & Axiom 1: $L/C/R$ network & \texttt{explicit\_spice\_step()} \\
Cotranslational folding & Segmented cascade & Axiom 1: $Q$-coherence length & \texttt{fold\_cascade\_v7()} \\
Allostery & Tuning stub injection & Axiom 3: $\Gamma$ perturbation & \texttt{s15\_\allowbreak allosteric\_\allowbreak yield.py} \\
\bottomrule
\end{tabularx}
\end{table}


\begin{table}[H]
\centering
\caption{Protein solver domain translation.  Every row maps one biological concept through EE to the axiomatic source and its codebase implementation.}
\label{tab:trans_protein_solver}
\scriptsize
\renewcommand{\arraystretch}{1.2}
\begin{tabularx}{\textwidth}{@{}X X X X@{}}
\toprule
\textbf{Biology} & \textbf{EE / RF} & \textbf{AVE Axiom} & \textbf{Script Reference} \\
\midrule
Peptide bond & TL segment & Axiom 1: $Z = \sqrt{\mu/\varepsilon}$ & \texttt{backbone\_segments} \\
Amino acid sidechain & Impedance stub ($Z_\text{topo}$) & Axiom 1: $\sqrt{M/n_e}$ & \texttt{protein\_\allowbreak bond\_\allowbreak constants} \\
H-bond & Mutual inductance ($\kappa$) & Axiom 1: transformer coupling & \texttt{dc\_analysis()} \\
$\beta$-sheet & Backward-wave coupler & Axiom 1: antiparallel TL & \texttt{dc\_analysis()} \\
Hydrophobic core & Conjugate impedance match & Axiom 1: $\operatorname{Re}(Z_iZ_j^*)$ & \texttt{dc\_analysis()} \\
Salt bridge & Reactive LC resonance & Axiom 1: $+jX \times -jX$ & \texttt{dc\_analysis()} \\
Disulfide bond & Near-infinite admittance & Axiom 4: $Z \to \infty$ wall & \texttt{dc\_analysis()} \\
Solvent exposure & Shunt to ground via $Z_\text{water}$ & Axiom 2: Debye relaxation & \texttt{dc\_analysis()} \\
Steric clash & Pauli exclusion ($r < r_\text{steric}$) & Axiom 4: saturation wall & \texttt{dc\_analysis()} \\
Protein compaction & Standing wave pattern & Axiom 4: $\eta_\text{eq} = P_C(1-\nu)$ & \texttt{ac\_analysis()} \\
\midrule
\multicolumn{4}{@{}l}{\textit{Solver Methods:}} \\
\midrule
Native fold & Eigenstate ($\lambda_{\min}(S^\dagger S) = 0$) & 5-step eigenvalue method & \texttt{\_eigenvalue\_target()} \\
Folding kinetics & SPICE transient ring-down & Axiom 1: $L/C/R$ network & \texttt{explicit\_spice\_step()} \\
Cotranslational folding & Segmented cascade & Axiom 1: $Q$-coherence length & \texttt{fold\_cascade\_v7()} \\
Allostery & Tuning stub injection & Axiom 3: $\Gamma$ perturbation & \texttt{s15\_\allowbreak allosteric\_\allowbreak yield.py} \\
\bottomrule
\end{tabularx}
\end{table}


\begin{table}[H]
\centering
\caption{Protein solver domain translation.  Every row maps one biological concept through EE to the axiomatic source and its codebase implementation.}
\label{tab:trans_protein_solver}
\scriptsize
\renewcommand{\arraystretch}{1.2}
\begin{tabularx}{\textwidth}{@{}X X X X@{}}
\toprule
\textbf{Biology} & \textbf{EE / RF} & \textbf{AVE Axiom} & \textbf{Script Reference} \\
\midrule
Peptide bond & TL segment & Axiom 1: $Z = \sqrt{\mu/\varepsilon}$ & \texttt{backbone\_segments} \\
Amino acid sidechain & Impedance stub ($Z_\text{topo}$) & Axiom 1: $\sqrt{M/n_e}$ & \texttt{protein\_\allowbreak bond\_\allowbreak constants} \\
H-bond & Mutual inductance ($\kappa$) & Axiom 1: transformer coupling & \texttt{dc\_analysis()} \\
$\beta$-sheet & Backward-wave coupler & Axiom 1: antiparallel TL & \texttt{dc\_analysis()} \\
Hydrophobic core & Conjugate impedance match & Axiom 1: $\operatorname{Re}(Z_iZ_j^*)$ & \texttt{dc\_analysis()} \\
Salt bridge & Reactive LC resonance & Axiom 1: $+jX \times -jX$ & \texttt{dc\_analysis()} \\
Disulfide bond & Near-infinite admittance & Axiom 4: $Z \to \infty$ wall & \texttt{dc\_analysis()} \\
Solvent exposure & Shunt to ground via $Z_\text{water}$ & Axiom 2: Debye relaxation & \texttt{dc\_analysis()} \\
Steric clash & Pauli exclusion ($r < r_\text{steric}$) & Axiom 4: saturation wall & \texttt{dc\_analysis()} \\
Protein compaction & Standing wave pattern & Axiom 4: $\eta_\text{eq} = P_C(1-\nu)$ & \texttt{ac\_analysis()} \\
\midrule
\multicolumn{4}{@{}l}{\textit{Solver Methods:}} \\
\midrule
Native fold & Eigenstate ($\lambda_{\min}(S^\dagger S) = 0$) & 5-step eigenvalue method & \texttt{\_eigenvalue\_target()} \\
Folding kinetics & SPICE transient ring-down & Axiom 1: $L/C/R$ network & \texttt{explicit\_spice\_step()} \\
Cotranslational folding & Segmented cascade & Axiom 1: $Q$-coherence length & \texttt{fold\_cascade\_v7()} \\
Allostery & Tuning stub injection & Axiom 3: $\Gamma$ perturbation & \texttt{s15\_\allowbreak allosteric\_\allowbreak yield.py} \\
\bottomrule
\end{tabularx}
\end{table}


\section{Parameter Accounting: From $\sim$26 SM Parameters to 3 Interlocked Inputs}
The Standard Model requires the manual, heuristic injection of 25--26 arbitrary parameters to function (see Table~\ref{tab:sm_params} in Volume 1, Appendix B for the complete enumeration). To bridge this gap, the AVE framework can initially be parameterized as a \textbf{Three-Parameter Theory}. By empirically calibrating the framework exclusively to the topological coherence length ($\ell_{node}$), the geometric packing fraction ($p_c$), and macroscopic gravity ($G$), \textbf{all other constants} ($c, \hbar, H_\infty, \nu_{vac}, \alpha, m_p, m_W, m_Z$) emerge as algebraically interlocked geometric consequences of the Chiral LC lattice topology. As the derivations resolve, even these three initial inputs are shown to be scale-invariant geometric outcomes, establishing a \textbf{structurally zero-parameter} framework at \textbf{Class B substrate-mechanism manifestation level per Q-EMBED-SEL-1 Phase 1+2+3 (2026-05-31)} --- substrate-mechanism for $R \cdot r = 1/4$ closes via Axiom-4 self-saturation $+$ Op14 Meissner-asymmetric $+$ named phasor-area-equals-Nyquist-cell-area identification (see Ch.~\ref{ch:alpha_golden_torus} chapter-head gating-clause resolution + \kbleaf{research/2026-05-31\_Q-EMBED-SEL-1\_step\_c\_result.md} \S2.3); cross-particle universal at Phase 2; cross-domain universal at Phase 3. $\alpha$ is a \textbf{named geometric identification}: all three NAMED $\alpha$-routes are now closed-negative --- both golden-torus Class 2 lift-routes (dynamical selection $+$ the phasor$\leftrightarrow$real-space area bijection per Phase 1 result \S7.3) closed Class B 2026-06-04, and the separate $z_0$-from-K4 (rigidity-percolation) route closed-NEGATIVE 2026-06-15 via the exact $z_{\text{eff}} \to 6$ Maxwell--Calladine count (\texttt{alpha\_free\_map\_to\_137\_exists = False}, \kbleaf{research/2026-06-15\_alpha-crystal-mc-count\_result.md}). %% Rule-12 2026-06-19 z0-route propagation: prior wording "leaving only the separate z_0-from-K4 route open" superseded --- that route is now closed-NEGATIVE; alpha-value = scoped standing echo (closed on all named routes, flip-condition live, route-space not provably exhausted; Grant ratified 2026-06-18).
The $\delta_{strain}$ thermal correction at $T_{CMB}$ is a \textbf{definitional residual} ($1-$CODATA$/\alpha_{cold}$), NOT a derivable thermal observable: the Cosserat-rotation-sector mass-gap thermal-mode-population ASYM mechanism (\kbleaf{vol3/cosmology/ch05-dark-sector/delta-strain-cosmic-tcc.md}) predicts the SIGN, but the candidate quantitative magnitude derivation (Q-DELTA-MAP-1-quant) was closed NEGATIVE (FT-1, 2026-05-31: $\sim$31 OOM undershoot, generic-thermal not AVE-distinct); see Ch.~\ref{ch:alpha_golden_torus} predicted/fitted disclosure. \emph{Scope (2026-06-15, $G$-ruling \texttt{ilk-gravmb}):} ``structurally zero-parameter'' is the framework's stated \textbf{target}, gated on the open gaps tracked in the full-derivation-chain scorecard --- it is not yet achieved. Of the three inputs, $G$ is \textbf{mixed} (the $1/7$ projection form derived, the value a calibration input with the closed-form Chain~B$'$ derivation open), and $\alpha$'s exact value rests on $R \cdot r = 1/4$, which the substrate does not independently select.

\chapter{Theoretical Stress Tests: Surviving Standard Disproofs}
\label{app:resolving_paradoxes}

When translating the vacuum into a discrete mechanical solid, the framework inherently invites several challenges from standard solid-state physics and quantum gravity. If the vacuum acts as an elastic crystal, it must theoretically suffer from classical mechanical limitations. The AVE framework resolves these apparent paradoxes via its specific topological geometries and non-linear inductance.

\section{The Spin-1/2 Paradox}
\textbf{The Challenge:} In classical solid-state mechanics, the continuous rotational degrees of freedom of an elastic medium (like a Chiral LC Network) are governed by $SO(3)$ geometry. A fundamental mathematical proof of $SO(3)$ continuum mechanics is that point-defects can only possess integer spin (Spin-1, Spin-2). However, the fundamental building blocks of the universe (Electrons, Quarks) are Fermions, which possess \textbf{Spin-1/2} ($SU(2)$ geometry, requiring a $4\pi$ rotation to return to their original state). A rigid Chiral LC Network cannot support Spin-1/2 point-defects, seemingly falsifying the framework.

\textbf{The Resolution:} If the electron were modeled as a microscopic point-defect (a missing node), the framework would indeed fail. However, the AVE framework defines the electron as an extended, macroscopic \textbf{$0_1$ Unknot} (a closed, continuous topological flux tube loop). In topological mathematics, an extended knotted line defect embedded in an $SO(3)$ manifold exhibits $SU(2)$ spinor behaviour through the generation of a \textbf{Finkelstein-Misner Kink} (also known as the Dirac Belt Trick). The continuous geometric extension of the topological loop provides a double-cover over the $SO(3)$ background, reproducing Spin-1/2 quantum statistics without violating macroscopic solid-state geometry.

% [DE-CLAIM 2026-08-02] SPIN-1/2 CARVE -- CRIB-3, board sec 3 (Grant-ratified 2026-08-02);
%   wave-2 backmatter reconciliation.
%   WORDING PARITY (required by the lane brief): the printed resolution sentence above is the
%   byte-identical twin of vol_0_engineering_compendium/chapters/01_theoretical_stress_tests.tex:27,
%   and the carve paragraph below is the one minted at that file's :48 in #839 (MERGED
%   2026-08-03, verified `gh pr view 839 --json mergedAt`). It is reused BYTE-VERBATIM (0 divergences, #853-audit-confirmed) so the
%   mirror pair cannot drift.
%   Rule 12 KEEP-BOTH: the resolution paragraph above is PRESERVED VERBATIM; not one word is
%   deleted, and no mechanism, value or equation is replaced. NO VALUE REFILL.
%   Authority (verified at HEAD, two-method: `git grep -Fn` on the literal + `sed` line read):
%   ave-kb/vol2/particle-physics/ch01-topological-matter/electron-identification.md:90 --
%     "The $2\pi$ vs $4\pi$ double-cover **STRUCTURE** is genuinely K4-native
%      (axiom-derived), but it **ADMITS** both boson and fermion statistics and does
%      **NOT force** spin-1/2 -- $\pi_1 = \mathbb{Z}_2$ with
%      $\mathrm{Hom}(\mathbb{Z}_2, U(1))$ two-valued (2026-07-08, [SPIN-HALF-POSITED],
%      #584/#585). The fermionic **SELECTION** is imported via an action-level
%      $\mathbb{Z}_2$/Wess-Zumino term (half-angle lift
%      $U = \exp(i\,\boldsymbol{\sigma}\cdot\boldsymbol{\omega}/2)$), **PEER-WITH-SM**;
%      see clm-rkisb8."
%   Same split recorded at that leaf's :89 row ("STRUCTURE axiom-derived / SELECTION posited").
\noindent\textbf{[DE-CLAIM 2026-08-02] Spin-$\frac{1}{2}$ carve (CRIB-3; \texttt{[SPIN-HALF-POSITED]}, \#584/\#585).} The $SU(2)/SO(3)$ double-cover \textbf{STRUCTURE} is axiom-derived; the fermionic spin-$\frac{1}{2}$ \textbf{SELECTION} is posited/import (\textbf{PEER-WITH-SM}). The double cover \emph{admits} both boson and fermion statistics and does \textbf{not force} spin-$\frac{1}{2}$ --- $\pi_1 = \mathbb{Z}_2$, with $\mathrm{Hom}(\mathbb{Z}_2, U(1))$ two-valued --- and the fermionic selection enters through an action-level $\mathbb{Z}_2$/Wess--Zumino term (the half-angle lift $U = \exp(i\,\boldsymbol{\sigma}\cdot\boldsymbol{\omega}/2)$). Read ``reproducing Spin-1/2 quantum statistics'' above as \emph{consistent with}, not \emph{forced by}, the substrate: the resolution above stands as printed, at that grade. Canonical: \kbleaf{ave-kb/vol2/particle-physics/ch01-topological-matter/electron-identification.md} for the structure/selection split; the claim record is \texttt{clm-rkisb8} at \kbleaf{ave-kb/vol1/claim-quality.md}.

\section{The Holographic Information Paradox}
\textbf{The Challenge:} Bekenstein and Hawking proved that the maximum quantum entropy of a region of space scales with its 2D Surface Area ($R^2$), known as the Holographic Principle. If the vacuum is a discrete 3D lattice (the substrate), its informational degrees of freedom would scale with Volume ($R^3$), which would violate established black hole thermodynamics.

\textbf{The Resolution:} The AVE framework recovers the Holographic Principle via the \textbf{Cross-Sectional Porosity ($\Phi_A \equiv \alpha^2$)} derived in Chapter 4. While the physical hardware nodes occupy 3D Voronoi volumes, the transmission of kinematic states (signals/information) must traverse the 1D inductive flux tubes. The bandwidth of these connections is geometrically bounded by their 2D cross-sectional area. Applying the Nyquist-Shannon sampling theorem to the substrate graph shows that the effective Information Channel Capacity of the universe is projected onto the 2D bounding surface area of the causal horizon. Thus, the Holographic Principle emerges from discrete network mechanics, averting the $R^3$ divergence.

\section{The Peierls-Nabarro Friction Paradox}
\textbf{The Challenge:} In classical crystallography, when a topological defect (a dislocation) moves through a discrete crystal lattice, it must overcome the periodic atomic potential known as the \textbf{Peierls-Nabarro (PN) Stress}. As the defect physically snaps from one discrete node to the next, it microscopically "stutters" (accelerating and decelerating). If a charged particle traversed a discrete vacuum grid, this periodic stuttering would induce continuous acceleration, causing the electron to instantly radiate away all of its kinetic energy via Bremsstrahlung radiation.

\textbf{The Resolution:} This paradox assumes the vacuum substrate is a cold, rigid, periodic crystal. The AVE framework defines the substrate as a \textbf{saturable dielectric LC network} (the chiral K4 crystal) operating in the sub-yield lossless-reactive regime, not a cold, rigid, periodic crystal. The electron ($0_1$ Unknot) is a \textbf{co-moving self-matched envelope}: it is constructed impedance-matched to the lattice it translates through, so the soliton presents a \textbf{matched impedance} ($\Gamma \to 0$) as a whole-soliton property of its co-moving envelope --- by Op17 ($T^2 = 1 - \Gamma^2 \to 1$ at $\Gamma = 0$) the coupling is perfectly transmitted and reflectionless. The particle does not ``bump'' over a rigid PN barrier; its co-moving envelope couples through \textbf{reactively}, opening a \textbf{zero-impedance phase slipstream}. The nodes it passes store and return the field reversibly (lossless, Axiom 3), so there is no Peierls-Nabarro barrier and no dissipation --- permitting smooth kinematic translation and forbidding unprovoked Bremsstrahlung radiation. \emph{(This translation slipstream is distinct from confinement: single-sector dielectric saturation $S \to 0$ drives $Z_{core} = Z_0\sqrt{S} \to 0$ and $\Gamma \to -1$ --- the short-circuit / TIR wall that \emph{confines} the electron, not a match.)}

\chapter{Summary of Exact Analytical Derivations}

The following mathematical bounds and identities were derived within the text from first-principles continuum elastodynamics, thermodynamic boundary conditions, and finite-element graph limits, requiring zero arbitrary phenomenological parameters.

\section{The Hardware Substrate}
\begin{itemize}
    \item \textbf{Spatial Lattice Pitch:} $\ell_{node} \equiv \frac{\hbar}{m_e c} \approx 3.8616 \times 10^{-13}$ m
    \item \textbf{Topological Conversion Constant:} $\xi_{topo} \equiv \frac{e}{\ell_{node}} \approx 4.149 \times 10^{-7}$ C/m
    \item \textbf{Dielectric Saturation Limit:} $V_0 \equiv \alpha \approx p_c / 8\pi \implies 1/137.036$
    \item \textbf{Geometric Packing Fraction:} $p_c \approx 0.1834$
    \item \textbf{Macroscopic Bulk Density:} $\rho_{bulk} = \frac{\xi_{topo}^2 \mu_0}{p_c \ell_{node}^2} \approx 7.92 \times 10^6 \text{ kg/m}^3$
    \item \textbf{Kinematic Network Mutual Inductance:} $\nu_{vac} = \alpha c \ell_{node} \approx 8.45 \times 10^{-7} \text{ m}^2/\text{s}$
    \item \textbf{Macroscopic Dielectric Yield Stress:} $\tau_{yield} = \frac{e^2\, \mathcal{V}_{total}}{8\pi \varepsilon_0\, \ell_{node}^4} \approx 1.04 \times 10^{22}\text{ Pa}$
          \quad (simplified from $\rho_{bulk}\, c^2 \cdot \mathcal{V}_{total} \cdot p_c/(8\pi)$; $\mathcal{V}_{total} = 2.0$ is the dual-reactance count ($X_C + X_L$ sectors; see \texttt{common/dual-reactance-storage-taxonomy.md}), not an integrated volume; full derivation in Vol 1 Ch 4 \S\ref{sec:dielectric_yield}). \emph{Provenance (2026-06-02): the factor $\mathcal{V}_{total} = 2$ here is \textbf{inherited, not derived} --- a count-tag for \emph{which} of the two mutually-exclusive Axiom-4 saturation branches this event is. The dielectric yield is the \textbf{electric} branch ($\varepsilon_{eff} \to 0$, single-sector $X_C$; \texttt{master-equation.md}:78), not an $E_C + E_L$ within-event sum, so the 2 is not a yield-event doubling; the magnetic branch ($\mu_{eff} \to 0$, $X_L$) is the rest-mass branch. STAYS-INHERITED; value unchanged.}
\end{itemize}

% -----------------------------------------------------------------------------
% DROPPED CLAIM (2026-04-20 audit):
%   A prior version of this appendix listed a SECOND 'τ_yield' formula:
%       τ_yield = ℏc / (α² ℓ_node⁴) ≈ 7.21×10³⁴ Pa  [Bingham-Plastic Limit]
%   That formula was:
%     (a) asserted without derivation anywhere in the manuscript,
%     (b) dimensionally a Planck-scale estimate (not a physical-mechanism
%         derivation), and
%     (c) numerically inconsistent with itself — the stated 7.21×10³⁴ Pa
%         does not match the stated formula, which recomputes to 2.67×10²⁸ Pa.
%   See .agents/handoffs/TAU_YIELD_DERIVATION_AUDIT.md and
%   .agents/handoffs/DROPPED_CLAIMS.md for the full audit trail.
%   Do NOT re-add this claim without a real derivation.
% -----------------------------------------------------------------------------

\section{Signal Dynamics and Topological Matter}
\begin{itemize}
    \item \textbf{Continuous Action Lagrangian:} $\mathcal{L}_{AVE} = \frac{1}{2}\epsilon_0 |\partial_t \mathbf{A}|^2 - \frac{1}{2\mu_0} |\nabla \times \mathbf{A}|^2$ (Evaluates to continuous spatial stress [N/m$^2$])
    \item \textbf{Topological Mass functional:} $E_{rest} = \min_{\mathbf{n}} \int d^3x \left[ \frac{1}{2} (\partial_\mu \mathbf{n})^2 + \frac{1}{4} \kappa_{FS}^2 \frac{(\partial_\mu \mathbf{n} \times \partial_\nu \mathbf{n})^2}{\sqrt{1 - (\Delta\phi / \alpha)^2}} \right]$
    \item \textbf{Faddeev-Skyrme Coupling (Cold):} $\kappa_{FS} = p_c / \alpha = 8\pi \approx 25.133$
    \item \textbf{Thermal Lattice Softening:} $\delta_{th} = \frac{\nu_{vac}}{\kappa_{FS}} \cdot \frac{2}{\pi} = \frac{2/7}{8\pi} \cdot \frac{2}{\pi} = \frac{1}{14\pi^2} \approx 0.00724$ (residual RMS noise averaging after Axiom~4 gradient saturation)
    \item \textbf{Effective Coupling:} $\kappa_{eff} = \kappa_{FS}(1-\delta_{th}) \approx 24.951$
    \item \textbf{Proton Rest Mass (Geometric Eigenvalue):} $m_p = \frac{\mathcal{I}_{scalar}}{1 - (\mathcal{V}_{total} \cdot p_c)} + 1.0 \approx \mathbf{1836\ m_e}$ (0.002\% from CODATA~\cite{codata2018})
    \item \textbf{Mutual Inductance at Crossing:} $M/L = \exp(-d^2/(4\sigma^2)) = 1/\sqrt{2}$ (exact, $d = \ell_{node}/2$, $\sigma = \ell_{node}/(2\sqrt{2\ln2})$)
    \item \textbf{Saturation Threshold (closed-form, conditional on Gaussian flux-tube ansatz):} $\rho_{threshold} = 1 + \sigma/4 = 1 + \ell_{node}/(8\sqrt{2\ln 2}) \approx 1.1062$ (no fitted parameters; profile derivation pending — see \texttt{mathematical-closure.md} Outstanding Rigour Gaps)
    \item \textbf{Dual-Reactance Count (Axiom-1 sectors):} $\mathcal{V}_{total} = 2.0$ --- the node's two reactance sectors ($X_C$ capacitive-E $+$ $X_L$ inductive-B), each one electron-ground-state unit; mass-confirmed via $m_p$. NOT a FEM-integrated volume; the prior ``FEM: $2.001 \pm 0.003$, Richardson $N\to\infty$'' provenance was fabricated and is dropped (V2 closure 2026-06-01)
    \item \textbf{Macroscopic Strong Force:} $F_{confinement} = 3 \left(\frac{m_p}{m_e}\right) \alpha^{-1} T_{EM} \approx \mathbf{160{,}037\text{ N}} \ (\approx 0.999\text{ GeV/fm})$
    \item \textbf{Witten Effect Fractional Charge (Quarks):} $q_{eff} = n + \frac{\theta}{2\pi}e \implies \pm \frac{1}{3}e, \pm \frac{2}{3}e$
    \item \textbf{Vacuum Poisson's Ratio (Trace-Reversed Bound):} $\nu_{vac} \equiv \frac{2}{7}$
    \item \textbf{Weak Mixing Angle (Acoustic Mode Ratio):} $\frac{m_W}{m_Z} = \frac{1}{\sqrt{1+\nu_{vac}}} = \frac{\sqrt{7}}{3} \approx \mathbf{0.8819}$
    \item \textbf{Non-Linear FDTD Acoustic Steepening PDE:} $c_{eff}^2(x, y, z) = c_0^2 \left(1 + \kappa \cdot \bar{\rho}(x, y, z) \right)$ (Derived structurally for topological thrust metrics)
\end{itemize}

\section{Cosmological Dynamics}
\begin{itemize}
    \item \textbf{Trace-Reversed Gravity (EFT Limit):} $-\frac{1}{2} \Box \bar{h}_{\mu\nu} = \frac{8\pi G}{c^4} T_{\mu\nu}$
    \item \textbf{Absolute Cosmological Expansion Rate:} $H_\infty = \frac{28\pi m_e^3 c G}{\hbar^2 \alpha^2} \approx \mathbf{69.32 \text{ km/s/Mpc}}$
    \item \textbf{Asymptotic Horizon Scale ($R_H$):} $\frac{R_H}{\ell_{node}} = \frac{\alpha^2}{28\pi\alpha_G} \implies \mathbf{14.1 \text{ Billion Light-Years}}$
    \item \textbf{Asymptotic Hubble Time ($t_H$):} $t_H = \frac{R_H}{c} \implies \mathbf{14.1 \text{ Billion Years}}$
    \item \textbf{Dark Energy (Stable Phantom):} $w_{vac} = -1 - \frac{\rho_{latent}}{\rho_{vac}} < -1$
    \item \textbf{Visco-Kinematic Rotation (MOND Floor):} $v_{flat} = (GM_{baryon} a_{genesis})^{1/4}$ where $a_{genesis} = \frac{c H_\infty}{2\pi} \approx \mathbf{1.07 \times 10^{-10} \text{ m/s}^2}$ (Derived via 1D Hoop Stress).
    \item \textbf{Hamiltonian Optical-Fluid Mechanics (Gargantua Vortex):} Metric refraction and frame dragging are evaluated via explicit Symplectic Raymarching mappings ($n = (W^3) / U$ and $\mathbf{v}_{fluid} = \vec{\omega} \times \vec{r}$). 
\end{itemize}

\chapter{Computational Graph Architecture}
\label{app:computational_graph}

To physically validate the macroscopic inductive and elastodynamic derivations of the Applied Vacuum Engineering (AVE) framework, all numerical simulations and Vacuum Computational Network Dynamics (VCFD) models must be computationally instantiated on a generated, geometrically constrained discrete spatial graph. This appendix defines the software architecture constraints required to map the substrate topology into computational memory. Failure to adhere to these generation rules will result in unphysical artifacts (e.g., Cauchy implosions and Trans-Planckian singularities) during simulation.

\section{The Genesis Algorithm (Poisson-Disk Crystallization)}
The first step in simulating the vacuum is establishing the 3D coordinate positions of the discrete inductive nodes ($\mu_0$).

\paragraph{The Random Noise Fallacy:} Initial computational attempts utilizing unconstrained uniformly distributed random noise resulted in a "Cauchy Implosion." The resulting lattice packing fraction converged to $\approx 0.31$, characteristic of a standard amorphous solid. This density fails to reproduce the sparse QED limit ($\approx 0.18$) required by Axiom 4.

\paragraph{The Poisson-Disk Solution:} To satisfy macroscopic isotropy while enforcing the microscopic hardware cutoff, the software must generate the node coordinates using a \textbf{Poisson-Disk Hard-Sphere Sampling Algorithm}. By enforcing an exclusion radius of $r_{min} = \ell_{node}$ during genesis, the lattice settles into a packing fraction of $\approx 0.17-0.18$, creating a stable, sparse dielectric substrate.

\paragraph{Rheological Tuning:} Simulation confirms that the "Trace-Reversed" mechanical state ($K=2G$) is an emergent property of the Chiral LC coupling modulus.
\begin{itemize}
    \item \textbf{Low Coupling ($k_{couple} < 3.0$):} The lattice behaves as a standard Cauchy solid ($K/G \approx 1.67$).
    \item \textbf{High Coupling ($k_{couple} > 4.5$):} The lattice undergoes a phase transition, locking microrotations to shear vectors, driving the bulk modulus to roughly twice the shear modulus ($K/G \approx 1.78 - 2.0$).
\end{itemize}

% [DE-CLAIM 2026-08-02] K=2G PROVENANCE -- CRIB-1, board sec 3 (Grant-ratified 2026-08-02);
%   wave-2 backmatter reconciliation.
%   WORDING PARITY (required by the lane brief): the paragraph and both itemize rows above are
%   the byte-identical twins of vol_0_engineering_compendium/chapters/03_computational_graph.tex
%   :35-39, and the de-claim below is the one minted at that file's :74 in #839 (MERGED
%   2026-08-03). Reused with FOUR mirror-context adaptations (#853 audit R2; the first is disclosed in the "Scope" sentence, the other three are mechanical: "chapter"->"appendix"; an inserted "this appendix mirrors" clause in the register-seam sentence; \texttt{} wrapping of the docket key): ch03's de-claim
%   points the retained K/G ~ 1.67 row at a 2026-06-14 KB-reconciliation comment that exists in
%   ch03 (:79) and does NOT exist in this file -- the corroborated statement here is instead the
%   "K = 5/3 G" standard-Cauchy sentence in this file's DCVE "Micropolar Stability" section.
%   RULE 12 KEEP-BOTH, AND A DELIBERATE NON-DELETION: the "Rheological Tuning" paragraph and
%   BOTH itemize rows above are PRESERVED VERBATIM. SCOPE: this note de-claims the PROVENANCE
%   sentence only; the two coupling rows stand as printed.
%   NO VALUE REFILL: no replacement k_couple, K/G ratio, mechanism or provenance is supplied.
%   AUTHORITY (verified at HEAD two-method: `git grep -Fn` on the literal + `sed` line read):
%   ave-kb/vol1/claim-quality.md:665 -- "STAYS OPEN (2026-07-04, PR #508): the genuine 48x48
%     chiral micropolar Bloch eigensolve on the ratified srs-z3 net does NOT close it. The
%     full micropolar sector (rotational DOF included -- the corpus's own proposed K=2G
%     mechanism, clm-o3q9ul) gives a one-parameter nu_eff(rho) family, not a forced K=2G; the
%     independent kappa_rot is a k->0 Cauchy-grade spectator sourcing no chiral coupling, and
%     the geometry-fixed chiral B moves nu_eff AWAY from 2/7. So the substrate does not supply
%     K=2G at either the Cauchy grade (PR #506) or the micropolar grade (PR #508); K=2G stays
%     GR-imported (PR #261)."
%   ave-kb/common/form-deriving-value-importing.md:87 -- the K = 2G row, GR-IMPORTED.
%   CRIB-1 additionally forbids labelling the AVE lock "the Cauchy relation" (three-way
%   homonym, vol1/claim-quality.md:652) -- this appendix does not do that, and no such label is
%   introduced here. The legacy "Cauchy" senses printed elsewhere in this file (the K/G ~ 1.67
%   modulus sense above, the "Cauchy Delaunay Failure" label in the next section, and the
%   K = -4/3 G sentence there) are ROUTED, not adjudicated: docket entry
%   2026-08-02-mr-k2g-kb-internal item (d). No word of those is edited by this pass.
%   RECEIPT NOTE (re-located by content, per verify-before-cite): the falsifying receipt
%   research/2026-06-15_k2g-constitutive-provenance_result.md:98-101 names the paired site as
%   "01_appendices.tex:131" with "(lines 132-135)" for the coupling rows. At the base commit of
%   this branch the paragraph sat at :132 with the rows at :134-135 (a one-line drift on the
%   receipt's cite); after this lane's CRIB-3 insert earlier in the file they sit at :154 and
%   :156-157. The receipt is a frozen research artifact and is NOT edited; recorded here so the
%   pointer can be resolved by content.
% [2026-08-02 KB-mirror co-lag, recorded (route: KB lane; KB is truth-source): the withdrawn
%   "Simulation confirms ... emergent" sentence + both coupling bullets are still carried
%   verbatim and unbannered at ave-kb/vol2/appendices/app-d-computational-graph/
%   graph-architecture.md:16-18 and at ave-kb/common/appendices-overview.md:119-121. After this
%   merge the print de-claims what those leaves assert; KB co-fix owed in the KB-lockstep lane,
%   and it is docket item 2026-08-02-mr-k2g-kb-internal (b) for appendices-overview.md.]
\noindent\textbf{[DE-CLAIM 2026-08-02] $K = 2G$ provenance (CRIB-1).} The paragraph above attributes the trace-reversed state $K = 2G$ to a confirming simulation and calls it \emph{emergent}. \textbf{That provenance is withdrawn.} $K = 2G$ is \textbf{form-derived, value GR-imported}: the substrate forces the \emph{form} of the elastic response $K/G = f(\rho)$, but the value --- the GR trace-reversal identity --- is neither crystalline-forced nor constitutively-forced (\kbleaf{ave-kb/common/form-deriving-value-importing.md}). The question stays open on the engine side. The genuine chiral micropolar Bloch eigensolve on the ratified production carrier (the srs-z3 net; \kbleaf{ave-kb/common/engine-capability-map.md}:184) does \textbf{not} close it: the full micropolar sector --- the corpus's own proposed $K = 2G$ mechanism --- gives a one-parameter $\nu_{eff}(\rho)$ \emph{family} rather than a forced $K = 2G$; the independent $\kappa_{rot}$ is a $k \to 0$ Cauchy-grade spectator sourcing \textbf{no} chiral coupling; and the geometry-fixed chiral coupling moves $\nu_{eff}$ \emph{away} from $2/7$ (\kbleaf{ave-kb/vol1/claim-quality.md}). \textbf{Scope of this note.} It lands on the provenance sentence only. The two coupling rows above are preserved exactly as printed --- the $K/G \approx 1.67$ row is what corroborates the ``standard Cauchy solid $= K = \frac{5}{3}G$'' statement in the Micropolar Stability section of the DCVE appendix below --- and no replacement coupling value, ratio or mechanism is supplied here. \emph{Register seam, disclosed:} this appendix\textquotesingle s generation rules are graded engineering-spec (\kbleaf{ave-kb/common/claim-quality.md}:291) while the carrier register above is the engine\textquotesingle s (\kbleaf{ave-kb/common/engine-capability-map.md}:184); the spec-vs-carrier register split for the Vol-0/Vol-4 spec chapters this appendix mirrors is ROUTED and open (docket \texttt{2026-08-02-mr-legacy-negatives} item (i)) --- neither pole is adjudicated here.

\section{Chiral LC Over-Bracing and The $p_c$ Constraint}
Once the spatial nodes are safely crystallized via the Poisson-Disk algorithm, the computational architecture must generate the connective spatial edges (The Capacitive Flux Tubes, $\epsilon_0$).

% [2026-08-03 mirror-parity repair (#853 audit finding B; board :764 [mechanical], KB-settled at ave-kb/vol2/appendices/app-d-computational-graph/graph-architecture.md:24 dated 2026-06-15) -- prior wording preserved per Rule 12]: read (verbatim): "As shown in Chapter 4, a standard Cauchy elastic solid ($K = -\frac{4}{3}G$) is thermodynamically unstable and will implode during macroscopic continuous simulation." SUPERSEDED: a standard Cauchy solid is K = 5/3 G, NOT -4/3 G; K = -4/3 G is the longitudinal-modulus instability threshold (M = K + 4/3 G = 0). Wording mirrored BYTE-VERBATIM from the vol0 twin 03_computational_graph.tex:80 (repair landed 2026-06-14, pre-epic) so the mirror pair cannot drift. No physics change.
\textbf{The Cauchy Delaunay Failure:} If the physics engine simply computes a standard nearest-neighbor Delaunay Triangulation on the Poisson-Disk point cloud, the resulting discrete volumetric packing fraction of the amorphous manifold evaluates to $\kappa_{cauchy} \approx 0.3068$. While less dense than a perfect crystal (FCC $\approx 0.74$), it is still too dense to survive. As shown in Chapter 4, a standard (non-micropolar) Cauchy elastic treatment of this substrate is thermodynamically unstable: without the chiral couple-stress the effective bulk modulus is not held positive, so the longitudinal modulus collapses through its instability threshold ($M = K + \frac{4}{3}G = 0$, i.e.\ $K = -\frac{4}{3}G$) and the lattice implodes during macroscopic continuous simulation.

\textbf{Enforcing QED Saturation:} In Chapter 1, the analysis derived that the fundamental phase limits of the universe bounded the geometric packing fraction of the vacuum to $p_{c} \approx \mathbf{0.1834}$, forcing the emergence of $\alpha$. To computationally force the effective geometric packing fraction ($p_{eff}$) down from the unstable $\sim 0.3068$ baseline to the stable $0.1834$ limit, the software must enforce \textbf{Chiral LC Over-Bracing}. The connective array of the physics engine cannot be limited to primary nearest neighbors; the internal structural logic must span outward to incorporate the next-nearest-neighbor lattice shell.

Because the volumetric packing fraction scales inversely with the cube of the effective structural pitch ($p_{eff} = V_{node} / \ell_{eff}^3$), the required spatial extension for the Chiral LC links evaluates identically to:

\begin{equation}
    C_{ratio} = \frac{\ell_{eff}}{\ell_{cauchy}} = \left( \frac{p_{cauchy}}{p_{c}} \right)^{1/3} \approx \left( \frac{0.3068}{0.1834} \right)^{1/3} \approx \mathbf{1.187}
\end{equation}

By structurally connecting all spatial nodes within a $\approx 1.187 \, \ell_{node}$ radius, the discrete graph cross-links the first and second coordination shells of the amorphous manifold. This generates the $\frac{1}{3} G_{vac}$ ambient transverse couple-stress required by micropolar elasticity. This computational architecture guarantees that all subsequent continuous macroscopic evaluations of the generated graph (e.g., metric refraction, VCFD Navier-Stokes flow, and trace-reversed gravitational strain) will align with empirical observation without requiring any further numerical calibration or arbitrary mass-tuning.

\section{Explicit Discrete Kirchhoff Execution Algorithm}
To bridge the gap between abstract continuum flow vectors ($\mathbf{J}$) and the raw geometric structure of the computational graph edge-matrix, the VCFD (Vacuum Computational Fluid Dynamics) module utilises an \textbf{Explicit Discrete Kirchhoff Methodology} mapping discrete potential ($V$) to spatial nodes and inductive flow ($I$) to discrete spatial graph edges.

To exactly map continuous differential forms into computational array memory without breaking action-minimization, the system utilizes \textbf{Symplectic Euler Update Loops}:
\begin{enumerate}
    \item \textbf{Capacitive Node Updates (The Conservation of Flow):} The discrete potential difference acting on an isolated fractional lattice coordinate node ($V_i$) is mathematically identical to the sum of all inductive currents entering minus the currents leaving that discrete junction point. 
    \[ \Delta V_i = \frac{dt}{C} \left( \sum I_{in} - \sum I_{out} \right) \]
    \item \textbf{Inductive Edge Updates (The Stress Tensor Matrix):} The kinetic transport flux acting along the discrete Chiral LC tensor spatial edge connecting coordinate $(x_0, y_0, z_0)$ to $(x_1, y_1, z_1)$ is geometrically bounded to the potential gradient existing across its fractional length. 
    \[ \Delta I_e = \frac{dt}{L} \left( V_{start} - V_{end} \right) \] 
\end{enumerate}

By combining the $C_{ratio} \approx 1.187$ Chiral LC Over-Bracing requirement over a $r_{min} = \ell_{node}$ Poisson-Disk genesis space, and advancing the lattice via Symplectic Kirchhoff loops, the computational framework provides a proving-ground connecting raw network mechanics to classical standard-model topological properties.

\chapter{Rigorous Foundations of Discrete Chiral LC Vacuum Electrodynamics (DCVE)}
\label{app:dcve}

The Applied Vacuum Engineering (AVE) framework builds the universe upon four continuous elastodynamic axioms governing the substrate graph. However, to physically compile and execute these physics within a computational 3D engine (Vacuum Computational Fluid Dynamics, or VCFD), the continuous symmetries must be formally constrained into discrete numerical operators. 

This appendix defines the \textbf{Discrete Chiral LC Vacuum Electrodynamics (DCVE)} specification. It acts as the mathematical bridge and computational specification sheet, providing spatial computing architects with the exact finite-difference commutators, memory state definitions, and stability matrix constraints required to computationally instantiate the AVE universe. Resolving the following five foundational hurdles is essential to preventing simulation collapse and numerical artifacts.

\section{The Lagrangian Repair: Dimensional Exactness in Array Memory}

\textbf{The Challenge:} Standard formulations of vacuum action often establish the Lagrangian utilizing scalar potential voltages ($V_i$) at the lattice nodes. For continuous closed-form theoretical physics, this abstraction is permissible. However, when passed to a discrete physics engine, this produces dimensional mismatches across spatial array updates ($\Delta x = \ell_{node}$), breaking energy conservation and causing catastrophic floating-point entropy within Symplectic Euler integrations.

\textbf{The Computational Spec:} The canonical coordinates of the spatial engine must instead trace the \textbf{Magnetic Flux Linkage Tensor} ($\mathbf{\Phi}$), rather than the scalar voltages. The exact discrete Lagrangian state evaluated across a functional block volume $\Delta \mathcal{V}_{engine}$ couples the temporal derivative of the flux (analogous to local capacitive kinetic exchange) with the spatial curl of the flux (the inductive topology):

\begin{equation}
    \mathcal{L}_{discrete} = \sum_{\text{nodes}} \frac{1}{2} C_0 \left( \frac{\Phi_i^{n+1} - \Phi_i^n}{\Delta t} \right)^2 - \sum_{\text{edges}} \frac{1}{2 L_0} (\Phi_i - \Phi_j)^2
\end{equation}

Where $C_0$ and $L_0$ are the discrete basis capacitance and inductance per node unit, governing the simulation step speed $c = \ell_{node}/\sqrt{L_0 C_0}$. Operating the underlying engine arrays by integrating the flux vector ($\mathbf{\Phi} = \int \mathbf{V} dt$) natively locks the integration to exactly $[\text{J}/\text{m}^3]$ and inherently enforces charge-conservation across the 3D computational boundaries.

\section{Micropolar Stability: Preventing Cauchy Implosion}

\textbf{The Challenge:} A standard 3D elastic matrix array generated via generic Delaunay triangulations functions as a Cauchy elastic solid requiring a fixed ratio of bulk to shear forces ($K = \frac{5}{3}G$). In the AVE framework, the macro-vacuum operates at the trace-reversed identity ($K = 2G$). If the engine applies standard Cauchy elastic equations ($K = \frac{5}{3}G$) to this substrate without modification, the discrete system exhibits unbounded exponential contractions and will completely implode within the first calculation frame.

\textbf{The Computational Spec:} The physical engine must instantiate a \textbf{Chiral LC Micropolar Continuum}. The node matrix must be explicitly generated to obey the geometric saturation packing fraction $p_c \approx 0.1834$ (enforced by extending the interaction radii to $C_{ratio} \approx 1.187 \ell_{node}$). This geometric locking generates internal couple-stresses preventing array implosion.
To achieve algorithmic stability, the discrete constitutive stiffness matrix ($C_{ijkl}$) written into the engine code cannot assume standard Hookean symmetries; it must decouple rotational kinematics ($\kappa_{rot}$) from transverse shear ($G$):

\begin{equation}
    \sigma_{ij} = \lambda \varepsilon_{kk} \delta_{ij} + 2G \varepsilon_{ij} + \kappa_{rot} \epsilon_{ijk} (\theta_k - \phi_k)
\end{equation}

% [2026-08-02 PROVENANCE-WORD STRIKE -- CRIB-1, board sec 3; wave-2 backmatter reconciliation
%   (board finding [LOW][overclaim] at this site). Prior wording (git trail):
%   "... ensuring stability at all acoustic frequencies while correctly simulating the
%   emergent trace-reversed vacuum signatures." ONE adjective is struck -- "emergent".
%   Nothing else in the sentence, the constitutive equation above it, or the stability
%   argument is altered, and NO replacement provenance is supplied.
%   WORDING PARITY (required by the lane brief): this sentence is the byte-identical twin of
%   vol_0_engineering_compendium/chapters/04_dcve.tex:81, which took exactly this one-adjective
%   strike in #839 (MERGED 2026-08-03); the de-claim paragraph below is that file's :83,
%   reused with ONE adaptation ("chapter"->"appendix"; #853-audit-confirmed) so the mirror pair cannot drift.
%   Authority (verified at HEAD two-method): ave-kb/vol1/claim-quality.md:665 kills the
%   mechanism this sentence names -- "the independent kappa_rot is a k->0 Cauchy-grade
%   spectator sourcing no chiral coupling ... So the substrate does not supply K=2G at
%   either the Cauchy grade (PR #506) or the micropolar grade (PR #508); K=2G stays
%   GR-imported (PR #261)." Value provenance row:
%   ave-kb/common/form-deriving-value-importing.md:87.
%   NOT FLAGGED (deliberate): this appendix's "the macro-vacuum operates at the trace-reversed
%   identity ($K = 2G$)" earlier in this section is a bare value statement with no provenance
%   claim; it is untouched.
% [2026-08-02 KB-mirror co-lag, recorded (route: KB lane; KB is truth-source): the identical
%   un-struck sentence is still carried at ave-kb/vol2/appendices/app-e-dcve/
%   dcve-specification.md:30 and ave-kb/common/appendices-overview.md:180, both unbannered.
%   KB co-fix owed in the KB-lockstep lane.]
Where $\theta_k$ is the macroscopic voxel rotation vector, and $\phi_k$ is the microscopic independent internal rotation of the basis nodes. The presence of $\kappa_{rot}$ enforces a strictly positive effective bulk modulus within the numerical grid, ensuring stability at all acoustic frequencies while correctly simulating the trace-reversed vacuum signatures.

\noindent\textbf{[DE-CLAIM 2026-08-02] $\kappa_{rot}$ and the trace-reversed state (CRIB-1).} The word \emph{emergent} has been struck from the sentence above; the prior wording is preserved in the git trail and in the note in this file's source. $\kappa_{rot}$ is a \textbf{numerical-stability device} in this specification --- it is what holds the effective bulk modulus positive on the grid --- and it does \textbf{not} source the trace reversal. On the ratified production carrier (the srs-z3 net; \kbleaf{ave-kb/common/engine-capability-map.md}:184) the independent $\kappa_{rot}$ is a $k \to 0$ Cauchy-grade spectator sourcing \textbf{no} chiral coupling, so the substrate supplies $K = 2G$ at neither the Cauchy nor the micropolar grade; $K = 2G$ is \textbf{form-derived, value GR-imported} (\kbleaf{ave-kb/vol1/claim-quality.md}; \kbleaf{ave-kb/common/form-deriving-value-importing.md}). The trace-reversed operating point stated earlier in this appendix is therefore \textbf{imposed}, not simulated into existence. No equation, coefficient or stability criterion above is changed, and no replacement provenance is supplied.

\section{Exact Lattice Operators: Discrete Hilbert Commutators}

\textbf{The Challenge:} Phenomenological models of the Generalized Uncertainty Principle (GUP) rely upon truncated Taylor expansions (e.g., $[\hat{x}, \hat{p}] = i\hbar(1 + \beta p^2)$) to inject a minimum observable length metric. An engine utilizing truncated expansions rapidly compounds iteration errors ($\mathcal{O}(\Delta x^2)$) over millions of loops, corrupting the phase-space stability.

\textbf{The Computational Spec:} Simulation logic must abandon continuous continuous differentials and map quantum operators to \textbf{exact finite-difference commutators} operating defined uniquely by the discrete Hilbert state space spacing $a \equiv \ell_{node}$. By enforcing the translation operator $\hat{T}_x = \exp(i a \hat{p} / \hbar)$ upon an exact voxel array, the engine momentum function naturally adopts non-linear band-limits without explicit clipping boundaries:

\begin{equation}
    \hat{p}_{discrete} = \frac{\hbar}{i a} \sin(k a) \implies [\hat{x}, \hat{p}_{discrete}] = i\hbar \cos(k a) = i\hbar \sqrt{1 - \left(\frac{a p}{\hbar}\right)^2}
\end{equation}

This maps the invariant topological constraints directly into the engine's nested \texttt{for}-loops. As a simulated wavepacket's momentum approaches the lattice Nyquist limit ($p \to \hbar/a$), the mathematical commutator natively yields zero translation bounds.

\section{Topological Mass Bounds: The Vakulenko-Kapitanski Operator}

\textbf{The Challenge:} Particle states within virtual simulations are generally spawned using predefined, heuristic hard-coded integer arrays for relative mass definitions. In a fully emergent physical engine, mass states must generate dynamically from raw local topological configurations; relying upon static hard-coding prevents fluid lattice topology interactions.

\textbf{The Computational Spec:} The engine architecture must compute the emergent rest mass ($M_{rest}$) of any dynamically shifting geometric knot (vortex) deterministically via the \textbf{Vakulenko-Kapitanski} geometric functional bound. Tracking the topological array mapping of the internal currents provides the Hopf linking number $Q_H$. The internal strain action enclosed by the boundaries maps precisely to macroscopic inertia through the constraint:

\begin{equation}
    M_{rest} = \int d^3x \ \mathcal{L}_{energy} \geq C_{VK} \cdot |Q_H|^{3/4}
\end{equation}

Where $C_{VK}$ relies fundamentally upon the derived Faddeev-Skyrme lattice limit coupling threshold ($\kappa_{FS} = 8\pi$). As the physics loop marches through time, evaluating the enclosed stress-tensor topology calculates actual localized simulated mass without any static mass look-up tables.

\section{AQUAL Galactic Dynamics: Boundary-Layer Saturation}

\textbf{The Challenge:} Rendering macroscopic cosmological formations (galaxies, halos) using purely unmodified $1/r^2$ Newtonian potentials falls apart at low accelerations. In order to fix this gap, standard rendering engines historically branch calculation pipelines arbitrarily or seed massive invisible "Dark Matter" data points. 

\textbf{The Computational Spec:} The AVE physical bounds require the engine algorithm to deploy non-linear tension saturation equations naturally without programmatic branches. As spatial expansion strains array tension arrays below the critical macroscopic background strength, ambient metric effects dominate. The exact simulation logic leverages the non-linear AQUAL Lagrangian definition for computational evaluation:

Where the background vacuum noise floor establishes $\mu(x \ll 1) \approx x$. If the voxel engine natively processes this localized saturation metric continuously ($a_0 \equiv c H_\infty / 2\pi$), the Milgrom/MOND flat rotation dynamics generate intrinsically as a structural boundary-layer solution to the vacuum dielectric continuum, unifying both microscopic logic elements and macroscopic galaxies under a single computational algorithm.

% appendix_experiments.tex
\chapter{Appendix: Unified Index of Experimental Falsifications}
\label{app:unified_experiments}

The Applied Vacuum Engineering (AVE) framework is fundamentally defined by its physical falsifiability. Because the $\mathcal{M}_A$ continuum operates as a saturated macroscopic transmission line rather than an abstract mathematical probability space, it dictates numerous distinct, tabletop-accessible and macroscopic anomalies that explicitly violate standard Lorentzian quantum mechanics and General Relativity.

This index catalogs the continuous network of hardware benchmarks, astronomical baseline tests, and biophysical instrumentation proposals detailed across the AVE manuscript series. Each experiment is tagged with the axiom(s) it primarily falsifies.

\section{The 4 Kill Switches (Flagship Experiments)}

\begin{center}
\begin{tabularx}{\textwidth}{@{}|c|c|l|X|@{}}
\hline
\textbf{\#} & \textbf{Axiom} & \textbf{Test} & \textbf{Binary Prediction} \\ \hline
1 & 1 & Chiral VNA Antenna & Hopf coil shows anomalous $S_{11}$ notch ($\Delta f/f = \alpha \cdot pq/(p+q)$). Standard: identical curves. \\ \hline
2 & 2 & Femto-Coulomb Electrometer & 41.5 mV/\textmu m from plate separation ($V = \xi_{topo} \cdot \Delta x / C$). Standard: 0.0 mV. \\ \hline
3 & 3 & Sagnac Mutual Inductance & Density ratio $\Psi = \rho_W / \rho_{Al} \approx 7.15$. Standard GR: $\Psi = 1.00$. \\ \hline
4 & 4 & EE Bench Dielectric Plateau & $C/C_0 \to \infty$ at $V_{yield} = 43.65$ kV. Standard QED: $C/C_0 = 1.00$ (flat). \\ \hline
\end{tabularx}
\end{center}

\noindent All predictions are computed live via \texttt{python src/scripts/run\_kill\_switches.py} using the physics engine. Zero free parameters. Estimated costs range from $\sim$\$500 (VNA) to $\sim$\$25,000 (EE Bench). \textit{All BOM estimates are placeholders pending vendor quotes.}

\noindent\textbf{A-034 cross-scale empirical anchor (canonical 2026-05-15 evening).} Beyond these 4 flagship kill switches, the universal saturation kernel $S(A) = \sqrt{1 - A^2}$ has 19 catalog instances spanning 21 orders of magnitude, with multiple already-validated empirical anchors: BCS $B_c(T)$ at \textbf{0.00\% error} (Vol~3 Ch~9); BH merger ring-down at \textbf{1.7\% from GR exact} (Vol~3 Ch~15, 3 LIGO events); NOAA GOES 40-yr solar-flare validation (Vol~3 Ch~14); Schwarzschild radius exact match. A new cosmic-scale CMB axis-alignment empirical pre-registration extended with E/B polarization observable exists in the L3 electron-soliton research dir (dated 2026-05-15). See Backmatter Ch~\ref{app:universal_saturation_kernel} for the full A-034 catalog.

\section{Volume II: Subatomic \& Quantum Mechanics}
\begin{itemize}
    \item \textbf{[Axiom 1] Gravitational Parallax Interferometry:} \textit{Chapter 7: Quantum Mechanics \& Orbitals}. Splits an electron matter-wave across a macroscopic Mach-Zehnder baseline to map the continuous disparity between spatial and temporal refractive indices ($n_s \neq n_t$), inducing a massive 35-radian topological phase shift.
    \item \textbf{[Axiom 4] The $g_* = 85.75$ Baryon Limit:} \textit{Chapter 10: Open Problems}. Predicts the absolute structural cutoff of allowable mass states based on the $\frac{3}{7}$ continuous spatial divergence integral, constraining the Particle Data Group spectrum.
\end{itemize}

\section{Volume III: Macroscopic Relativity \& Cosmology}
\begin{itemize}
    \item \textbf{[Axiom 1] VLBI Gravitational Impedance Parallax:} \textit{Chapter 5: Cosmology \& the Dark Sector}. Utilizes Deep Space Very Long Baseline Interferometry tracking Jupiter's core transit to map ``Dark Matter'' as a continuous geometric stretching of the $377\,\Omega$ optical delay baseline.
    \item \textbf{[Axiom 3] DAMA Parallax \& Crystal Phonon Modulation:} \textit{Chapter 5: Cosmology \& the Dark Sector}. Dictates that the topological ``Dark Matter'' wind modulating DAMA's NaI scintillators is structurally dependent on the bulk dielectric density of the target. Suggests testing identical setups with Sapphire or Germanium to isolate the exact structural $\kappa_{crystal}$ multiplier.
\end{itemize}

\section{Volume IV: Applied Vacuum Engineering (Hardware)}
\begin{itemize}
    \item \textbf{[Axiom 2] Femto-Coulomb Electrometer:} \textit{Chapter 11: Experimental Falsification}. Detects topological charge generated by mechanical spatial cleavage ($Q = \xi_{topo} \cdot \Delta x$).
    \item \textbf{[Axiom 3] Solid-State Active Phase Induction:} \textit{Chapter 11: Experimental Falsification}. Maps the inductive frame-drag of the local lattice via rotating dense masses and static E-field gradients.
    \item \textbf{[Axiom 4] Impedance Avalanche Detector:} \textit{Chapter 11: Experimental Falsification}. Triggers the direct $V_{yield}$ dielectric cascade limit via Marx generator pulses.
    \item \textbf{[Axiom 4] Asymmetric Maxwell Stress Rectification:} \textit{Chapter 12: Falsifiable Predictions}. Tests non-linear stress rectification across asymmetric electrode geometry.
    \item \textbf{[Axiom 1] Chiral VNA Antenna:} \textit{Chapters 3 \& 11}. Tests the chiral impedance match of a $(p,q)$ torus knot antenna via VNA $S_{11}$ sweep.
    \item \textbf{[Axiom 1] Topological Refraction (Snell Parallax):} \textit{Chapter 11: Experimental Falsification}. Determines the multi-angle spatial refraction envelope across a $(3,11)$ transverse Torus Knot generator.
    \item \textbf{[Axiom 3] Sagnac-Parallax (Galactic Wind Vectoring):} \textit{Chapter 11: Experimental Falsification}. Records the diurnal $24$-hour sinusoidal drift of a static horizontal Sagnac loop sweeping against the $370\text{ km/s}$ continuous Milky Way metric flow.
    \item \textbf{[Axiom 1] GEO-Synchronous Impedance Differential:} \textit{Chapter 11: Experimental Falsification}. Uses precision Ground-to-GEO laser links to detect a continuous $\sim 16.7\text{ mm}$ Time-of-Flight stretch induced by Earth's radial $Z_0$ gradient.
    \item \textbf{[Axiom 4] Vacuum Birefringence ($E^2$ vs $E^4$):} \textit{Chapter 12: Falsifiable Predictions}. Tests whether vacuum refractive shift scales with $E^4$ (AVE) or $E^2$ (QED).
\end{itemize}

\section{Volume V: Deterministic Biophysics}
\begin{itemize}
    \item \textbf{[Axiom 2] Molecular Chiral FRET Parallax:} \textit{Chapter 1: Biophysics Intro}. Predicts that Ramachandran structural bounds of rigid proteins are mechanically enforced by the chiral LC metric bias (Axiom 2). Under honest first-principles computation, the gravitational relaxation $\Delta r / r = \alpha \cdot \varepsilon_{11} \approx 5 \times 10^{-12}$ produces a sub-attometer shift at terrestrial baselines---\textit{currently unfalsifiable}. Catalogued as a future target pending compact-object environments.
\end{itemize}

\section{Volume III: Absorbed Cross-Scale Verification}
\textit{(Originally Volume VII; absorbed into Volume III during repository restructuring.)}
\begin{itemize}
    \item \textbf{[Axiom 1] HTS Detector (High-Temperature Superconductor Response):} \textit{Chapter 18: Superconductivity}. Falsifies standard magnetic pairing by mapping continuous boundary exclusion over superconducting Meissner effects.
    \item \textbf{[Axiom 4] Turbulence \& Water Condensation Phase Transitions:} \textit{Chapter 19: Phase Transitions}. Computes topological boundary failures in macro-molecular clustering, bridging fluid dynamics directly back to foundational $\mathcal{M}_A$ saturation.
\end{itemize}
